\documentclass[../main.tex]{subfiles}
\begin{document}

\chapter{Scheduling}
\lessoninfo{
  video={P3L1},
  textbook={OSTEP \cite{ostep} ch.~7--10; Silberschatz et al.\ \cite{silberschatz2018} ch.~5; Fedorova et al.\ \cite{fedorova2004}},
  keywords={runqueue, FCFS, SJF, priority inversion, round robin, timeslice, MLFQ, O(1) scheduler, CFS, cache affinity, NUMA, simultaneous multithreading, hardware performance counter},
}

Lesson~2 introduced the CPU scheduler, which selects the next process to
run from the ready queue, and showed that the length of the timeslice trades
the overhead of scheduling against the time that ready processes wait.  This
lesson studies how the selection is made: scheduling algorithms from simple
run-to-completion policies to preemptive, priority-based and timesharing
schedulers, the runqueue data structures that implement them, and the two
schedulers of the Linux kernel.  It then turns to the questions raised by
multiprocessors and by hyperthreaded processors, following the paper of
Fedorova et al.\ \cite{fedorova2004}.  The performance metrics of Lesson~6
serve to compare the algorithms.

\section{Overview of CPU scheduling}
\label{sec:l07-overview}

The CPU scheduler decides how and when processes, and the threads within
them, access the shared CPUs.\source{P3L1, visual metaphor, scheduling
overview; OSTEP ch.~7; Silberschatz et al.\ ch.~5.}  It schedules tasks that
run user-level processes and threads as well as the kernel's own
kernel-level threads.  Following the Linux terminology of Lesson~5, this
lesson calls every such schedulable entity a \emph{task}.

Even the simplest scheduler embodies a goal.  Three elementary strategies
show the range of possible goals:
\begin{itemize}
  \item dispatch tasks immediately, in the order in which they arrive: the
    scheduling decision is trivial and its overhead is low;
  \item dispatch simple, short tasks first: this finishes the largest number
    of tasks early, and so maximizes the throughput, the number of tasks
    completed per unit time, over any interval shorter than the whole
    workload;
  \item dispatch complex, long tasks first: this keeps the resources
    occupied and maximizes the utilization of the CPU, the devices and
    memory.
\end{itemize}

The scheduler runs whenever the CPU may have to be given to a different
task:
\begin{itemize}
  \item the CPU becomes idle, for example because the running task blocks
    waiting for I/O;
  \item a new task becomes ready, for example a newly created task or a task
    whose I/O operation has completed; the scheduler must determine whether
    the new task is more important than the running one and should interrupt
    it;
  \item the timeslice of the running task expires.\margin{A timer tick
    notices the expiry: Linux kernels are built with \texttt{CONFIG\_HZ} of
    100, 250 or 1000, a tick every 10, 4 or \SI{1}{\milli\second}.}
\end{itemize}
The scheduler then selects one of the ready tasks and dispatches it: it
performs a context switch to the selected task, enters user mode, and sets
the program counter to the next instruction of the task, which continues to
execute from there.

Which task is selected is determined by the scheduling policy, or
algorithm.  How the selection is carried out depends on the data structure
that holds the ready tasks.

\begin{definition}[Runqueue]
  The data structure in which the scheduler keeps the tasks that are ready
  to run is its \term{runqueue}[ランキュー].
\end{definition}

The ready queue of Lesson~2 is such a runqueue.  Despite the name, a
runqueue need not be a FIFO queue: many schedulers use trees, or several
queues.  The runqueue and the
scheduling algorithm are tightly coupled.  An algorithm can be implemented
efficiently only with a runqueue that makes its selection cheap, and a
runqueue is designed for the algorithm that it serves.

\Needspace{7\baselineskip}
\section{Run-to-completion scheduling}
\label{sec:l07-rtc}

The first schedulers to consider never take the CPU away from a
task.\source{P3L1, run to completion scheduling, SJF performance; OSTEP
ch.~7; Silberschatz et al.\ ch.~5.}  We analyze them under simplifying
assumptions: a group of tasks is to be scheduled, the execution time of
every task is known in advance, there is no preemption, and there is a
single CPU.  The algorithms are compared by throughput, average completion
time, average wait time and CPU utilization (Lesson~6).  The completion time
of a task is measured from its arrival until it completes, and its wait
time from its arrival until it starts to run.

\begin{definition}[Run-to-completion scheduling]
  Under \term{run-to-completion scheduling}[実行完了型スケジューリング] a
  task, once dispatched, runs on the CPU until it completes; it is never
  preempted.
\end{definition}

\begin{definition}[FCFS]
  A scheduler that follows
  \term{first-come first-served scheduling}[先着順スケジューリング]%
  \index{FCFS|see{first-come first-served scheduling}} (FCFS) runs the tasks
  in the order in which they arrive.
\end{definition}

The runqueue of FCFS is a FIFO queue: an arriving task is appended at the
tail, and the scheduler takes the task at the head.  Both operations take
constant time.

\begin{definition}[SJF]
  A scheduler that follows the policy
  \term{shortest job first}[最短ジョブ優先]%
  \index{SJF|see{shortest job first}} (SJF) runs the tasks in the order of
  their execution times, shortest first.
\end{definition}

The runqueue of SJF must be ordered by execution time.  An ordered queue
lets the scheduler take the first task in constant time, but inserting a
task requires a search for its position, which takes time linear in the
length of the queue.  A balanced tree ordered by execution time inserts a
task in logarithmic time, and the scheduler takes its leftmost node.

\begin{figure}[htbp]
  \centering
  % FCFS, SJF and round robin (timeslice 1 s) for three tasks that arrive
% at time 0 in the order T1, T2, T3 with execution times 5 s, 2 s, 1 s.
% Usage: % FCFS, SJF and round robin (timeslice 1 s) for three tasks that arrive
% at time 0 in the order T1, T2, T3 with execution times 5 s, 2 s, 1 s.
% Usage: % FCFS, SJF and round robin (timeslice 1 s) for three tasks that arrive
% at time 0 in the order T1, T2, T3 with execution times 5 s, 2 s, 1 s.
% Usage: \input{figures/l07-gantt}   (width ~106 mm)
\begin{tikzpicture}[x=10mm, y=1mm, font=\footnotesize\sffamily,
  t1/.style={draw, fill=ln-accent!25},
  t2/.style={draw, fill=ln-box},
  t3/.style={draw, fill=white},
  lab/.style={anchor=east, align=right}]
  % FCFS
  \node[lab] at (-0.2,25) {FCFS};
  \draw[t1] (0,22) rectangle (5,28); \node at (2.5,25) {T1};
  \draw[t2] (5,22) rectangle (7,28); \node at (6,25) {T2};
  \draw[t3] (7,22) rectangle (8,28); \node at (7.5,25) {T3};
  % SJF
  \node[lab] at (-0.2,14) {SJF};
  \draw[t3] (0,11) rectangle (1,17); \node at (0.5,14) {T3};
  \draw[t2] (1,11) rectangle (3,17); \node at (2,14) {T2};
  \draw[t1] (3,11) rectangle (8,17); \node at (5.5,14) {T1};
  % round robin, timeslice 1 s
  \node[lab] at (-0.2,3) {\iflangja{RR（1\,s）}{RR (1\,s)}};
  \draw[t1] (0,0) rectangle (1,6); \node at (0.5,3) {T1};
  \draw[t2] (1,0) rectangle (2,6); \node at (1.5,3) {T2};
  \draw[t3] (2,0) rectangle (3,6); \node at (2.5,3) {T3};
  \draw[t1] (3,0) rectangle (4,6); \node at (3.5,3) {T1};
  \draw[t2] (4,0) rectangle (5,6); \node at (4.5,3) {T2};
  \draw[t1] (5,0) rectangle (8,6); \node at (6.5,3) {T1};
  \draw[densely dotted] (6,0) -- (6,6) (7,0) -- (7,6);
  % time axis
  \draw[->] (0,-3) -- (8.6,-3);
  \node[anchor=west] at (8.6,-3) {$t$\,(s)};
  \foreach \x in {0,...,8} {
    \draw (\x,-3) -- (\x,-4.2);
    \node[below, font=\scriptsize\sffamily] at (\x,-4.2) {\x};
  }
\end{tikzpicture}
   (width ~106 mm)
\begin{tikzpicture}[x=10mm, y=1mm, font=\footnotesize\sffamily,
  t1/.style={draw, fill=ln-accent!25},
  t2/.style={draw, fill=ln-box},
  t3/.style={draw, fill=white},
  lab/.style={anchor=east, align=right}]
  % FCFS
  \node[lab] at (-0.2,25) {FCFS};
  \draw[t1] (0,22) rectangle (5,28); \node at (2.5,25) {T1};
  \draw[t2] (5,22) rectangle (7,28); \node at (6,25) {T2};
  \draw[t3] (7,22) rectangle (8,28); \node at (7.5,25) {T3};
  % SJF
  \node[lab] at (-0.2,14) {SJF};
  \draw[t3] (0,11) rectangle (1,17); \node at (0.5,14) {T3};
  \draw[t2] (1,11) rectangle (3,17); \node at (2,14) {T2};
  \draw[t1] (3,11) rectangle (8,17); \node at (5.5,14) {T1};
  % round robin, timeslice 1 s
  \node[lab] at (-0.2,3) {\iflangja{RR（1\,s）}{RR (1\,s)}};
  \draw[t1] (0,0) rectangle (1,6); \node at (0.5,3) {T1};
  \draw[t2] (1,0) rectangle (2,6); \node at (1.5,3) {T2};
  \draw[t3] (2,0) rectangle (3,6); \node at (2.5,3) {T3};
  \draw[t1] (3,0) rectangle (4,6); \node at (3.5,3) {T1};
  \draw[t2] (4,0) rectangle (5,6); \node at (4.5,3) {T2};
  \draw[t1] (5,0) rectangle (8,6); \node at (6.5,3) {T1};
  \draw[densely dotted] (6,0) -- (6,6) (7,0) -- (7,6);
  % time axis
  \draw[->] (0,-3) -- (8.6,-3);
  \node[anchor=west] at (8.6,-3) {$t$\,(s)};
  \foreach \x in {0,...,8} {
    \draw (\x,-3) -- (\x,-4.2);
    \node[below, font=\scriptsize\sffamily] at (\x,-4.2) {\x};
  }
\end{tikzpicture}
   (width ~106 mm)
\begin{tikzpicture}[x=10mm, y=1mm, font=\footnotesize\sffamily,
  t1/.style={draw, fill=ln-accent!25},
  t2/.style={draw, fill=ln-box},
  t3/.style={draw, fill=white},
  lab/.style={anchor=east, align=right}]
  % FCFS
  \node[lab] at (-0.2,25) {FCFS};
  \draw[t1] (0,22) rectangle (5,28); \node at (2.5,25) {T1};
  \draw[t2] (5,22) rectangle (7,28); \node at (6,25) {T2};
  \draw[t3] (7,22) rectangle (8,28); \node at (7.5,25) {T3};
  % SJF
  \node[lab] at (-0.2,14) {SJF};
  \draw[t3] (0,11) rectangle (1,17); \node at (0.5,14) {T3};
  \draw[t2] (1,11) rectangle (3,17); \node at (2,14) {T2};
  \draw[t1] (3,11) rectangle (8,17); \node at (5.5,14) {T1};
  % round robin, timeslice 1 s
  \node[lab] at (-0.2,3) {\iflangja{RR（1\,s）}{RR (1\,s)}};
  \draw[t1] (0,0) rectangle (1,6); \node at (0.5,3) {T1};
  \draw[t2] (1,0) rectangle (2,6); \node at (1.5,3) {T2};
  \draw[t3] (2,0) rectangle (3,6); \node at (2.5,3) {T3};
  \draw[t1] (3,0) rectangle (4,6); \node at (3.5,3) {T1};
  \draw[t2] (4,0) rectangle (5,6); \node at (4.5,3) {T2};
  \draw[t1] (5,0) rectangle (8,6); \node at (6.5,3) {T1};
  \draw[densely dotted] (6,0) -- (6,6) (7,0) -- (7,6);
  % time axis
  \draw[->] (0,-3) -- (8.6,-3);
  \node[anchor=west] at (8.6,-3) {$t$\,(s)};
  \foreach \x in {0,...,8} {
    \draw (\x,-3) -- (\x,-4.2);
    \node[below, font=\scriptsize\sffamily] at (\x,-4.2) {\x};
  }
\end{tikzpicture}

  \caption{Three tasks that arrive at time 0 in the order T1, T2, T3, with
    execution times of \SI{5}{\second}, \SI{2}{\second} and
    \SI{1}{\second}, under FCFS, SJF and round robin with a timeslice of
    \SI{1}{\second} (\cref{sec:l07-rr}).  Dotted lines separate consecutive
    timeslices of the same task.}
  \label{fig:l07-gantt}
\end{figure}

\begin{example}
  \label{ex:l07-rtc}
  Three tasks arrive at time 0 in the order T1, T2, T3, with execution
  times of \SI{5}{\second}, \SI{2}{\second} and \SI{1}{\second}
  (\cref{fig:l07-gantt}).  Under FCFS they complete at 5, 7 and
  \SI{8}{\second}.  The throughput is $3/8 = 0.375$ tasks per second, the
  average completion time is $(5+7+8)/3 \approx \SI{6.67}{\second}$, and
  the average wait time is $(0+5+7)/3 = \SI{4}{\second}$.  SJF runs T3, T2,
  T1, which complete at 1, 3 and \SI{8}{\second}.  The throughput is still
  0.375 tasks per second, but the average completion time falls to
  $(1+3+8)/3 = \SI{4}{\second}$ and the average wait time to
  $(0+1+3)/3 \approx \SI{1.33}{\second}$.
\end{example}

The total work is the same in both orders, so the throughput measured over
the whole batch does not change.  By running short tasks first, however, SJF
has more tasks finished at any instant before the end of the batch, and it
lowers both averages.  In fact, for tasks that
arrive together and run to completion, no other order yields a smaller
average wait time.

\Needspace{7\baselineskip}
\section{Preemptive scheduling}
\label{sec:l07-preemptive}

In practice tasks arrive at different times, and a task that arrives may
deserve the CPU more than the task that is running.\source{P3L1, preemptive
scheduling: SJF + preempt, priority; OSTEP ch.~7--8; Silberschatz et al.\
ch.~5.}  It may be shorter, or more important.  A run-to-completion
scheduler makes it wait until the running task completes, however long that
takes.

\begin{definition}[Preemptive scheduling]
  Under \term{preemptive scheduling}[プリエンプティブスケジューリング] the
  scheduler may interrupt a running task before it completes and give the
  CPU to another task.
\end{definition}

SJF becomes preemptive as follows.  Whenever a task arrives, the scheduler
compares its execution time with the remaining execution time of the
running task.  If the new task is shorter, the running task is preempted,
returned to the runqueue, and the new task runs.\margin{This policy is also
known as \emph{shortest remaining time first}; OSTEP calls it
\emph{shortest time-to-completion first} (STCF).}

The execution time of a task is rarely known in advance.  The scheduler
therefore estimates it from the history of the task, using heuristics such
as the time the task ran the last time it was on the CPU, or a windowed
average of the times it ran the last $n$ times.\margin{Silberschatz et al.\
use an exponential average, $\tau_{n+1} = \alpha t_n + (1-\alpha)\tau_n$:
with $\alpha = 1/2$ only $\tau_n$ need be kept, and every older burst counts
half as much as the one after it.}

\subsection{Priority scheduling}

Tasks also differ in importance.  Kernel-level threads that perform
critical work for the operating system may be more important than threads
of user processes, and among user processes some are more urgent than
others.

\begin{definition}[Priority]
  The \term{priority}[優先度] of a task is a value that expresses its
  importance to the scheduler.  A scheduler that uses
  \term{priority scheduling}[優先度スケジューリング] runs the ready task of
  highest priority next, and preempts the running task when a task of higher
  priority becomes ready.
\end{definition}

The runqueue of a priority scheduler can consist of one queue per priority
level, from which the scheduler selects the first task of the highest
non-empty level; or it can be a tree ordered by priority.

\begin{example}
  \label{ex:l07-preemptive}
  Three tasks arrive as follows (\cref{fig:l07-preemptive}):

  \begin{center}\small
  \begin{tabular}{@{}lccc@{}}
    \toprule
    Task & Arrival (s) & Execution time (s) & Priority \\
    \midrule
    T1 & 0 & 7 & medium \\
    T2 & 1 & 2 & low \\
    T3 & 2 & 3 & high \\
    \bottomrule
  \end{tabular}
  \end{center}

  With preemptive SJF, T1 runs until T2 arrives at \SI{1}{\second}; T2 is
  shorter than the \SI{6}{\second} that remain of T1 and preempts it.  When
  T3 arrives at \SI{2}{\second}, only \SI{1}{\second} of T2 remains, so T2
  continues and completes at \SI{3}{\second}.  T3 is shorter than the rest
  of T1 and runs next, completing at \SI{6}{\second}; T1 completes at
  \SI{12}{\second}.  Measured from their arrivals, the tasks take 12, 2 and
  \SI{4}{\second}, \SI{6}{\second} on average.  With preemptive priority
  scheduling, T2 does not preempt T1, but T3 does: T3 runs from 2 to
  \SI{5}{\second}, T1 then completes at \SI{10}{\second}, and T2 only at
  \SI{12}{\second}.
\end{example}

\begin{figure}[htbp]
  \centering
  % Preemptive SJF and preemptive priority scheduling of three tasks:
% T1 arrives at 0 (7 s, medium), T2 at 1 (2 s, low), T3 at 2 (3 s, high).
% Usage: % Preemptive SJF and preemptive priority scheduling of three tasks:
% T1 arrives at 0 (7 s, medium), T2 at 1 (2 s, low), T3 at 2 (3 s, high).
% Usage: % Preemptive SJF and preemptive priority scheduling of three tasks:
% T1 arrives at 0 (7 s, medium), T2 at 1 (2 s, low), T3 at 2 (3 s, high).
% Usage: \input{figures/l07-preemptive}   (width ~100 mm)
\begin{tikzpicture}[x=7.5mm, y=1mm, font=\footnotesize\sffamily,
  >={Stealth[length=1.5mm]},
  t1/.style={draw, fill=ln-accent!25},
  t2/.style={draw, fill=ln-box},
  t3/.style={draw, fill=white},
  ttl/.style={anchor=south west, font=\footnotesize\sffamily\bfseries, inner sep=1pt}]
  % preemptive SJF
  \node[ttl] at (0,26.5) {\iflangja{(a) プリエンプティブ SJF}{(a) preemptive SJF}};
  \draw[t1] (0,20) rectangle (1,26);  \node at (0.5,23) {T1};
  \draw[t2] (1,20) rectangle (3,26);  \node at (2,23) {T2};
  \draw[t3] (3,20) rectangle (6,26);  \node at (4.5,23) {T3};
  \draw[t1] (6,20) rectangle (12,26); \node at (9,23) {T1};
  % priority
  \node[ttl] at (0,10.5) {\iflangja{(b) プリエンプティブ優先度スケジューリング}{(b) preemptive priority scheduling}};
  \draw[t1] (0,4) rectangle (2,10);   \node at (1,7) {T1};
  \draw[t3] (2,4) rectangle (5,10);   \node at (3.5,7) {T3};
  \draw[t1] (5,4) rectangle (10,10);  \node at (7.5,7) {T1};
  \draw[t2] (10,4) rectangle (12,10); \node at (11,7) {T2};
  % time axis
  \draw[->] (0,0) -- (12.8,0);
  \node[anchor=west] at (12.8,0) {$t$\,(s)};
  \foreach \x in {0,...,12} {
    \draw (\x,0) -- (\x,-1.2);
    \node[below, font=\scriptsize\sffamily] at (\x,-1.2) {\x};
  }
  % arrivals
  \node[anchor=east, font=\scriptsize\sffamily] at (-0.2,-9) {\iflangja{到着}{arrival}};
  \foreach \x/\t in {0/T1, 1/T2, 2/T3} {
    \draw[->] (\x,-10.5) -- (\x,-5.5);
    \node[below, font=\scriptsize\sffamily, inner sep=1pt] at (\x,-10.5) {\t};
  }
\end{tikzpicture}
   (width ~100 mm)
\begin{tikzpicture}[x=7.5mm, y=1mm, font=\footnotesize\sffamily,
  >={Stealth[length=1.5mm]},
  t1/.style={draw, fill=ln-accent!25},
  t2/.style={draw, fill=ln-box},
  t3/.style={draw, fill=white},
  ttl/.style={anchor=south west, font=\footnotesize\sffamily\bfseries, inner sep=1pt}]
  % preemptive SJF
  \node[ttl] at (0,26.5) {\iflangja{(a) プリエンプティブ SJF}{(a) preemptive SJF}};
  \draw[t1] (0,20) rectangle (1,26);  \node at (0.5,23) {T1};
  \draw[t2] (1,20) rectangle (3,26);  \node at (2,23) {T2};
  \draw[t3] (3,20) rectangle (6,26);  \node at (4.5,23) {T3};
  \draw[t1] (6,20) rectangle (12,26); \node at (9,23) {T1};
  % priority
  \node[ttl] at (0,10.5) {\iflangja{(b) プリエンプティブ優先度スケジューリング}{(b) preemptive priority scheduling}};
  \draw[t1] (0,4) rectangle (2,10);   \node at (1,7) {T1};
  \draw[t3] (2,4) rectangle (5,10);   \node at (3.5,7) {T3};
  \draw[t1] (5,4) rectangle (10,10);  \node at (7.5,7) {T1};
  \draw[t2] (10,4) rectangle (12,10); \node at (11,7) {T2};
  % time axis
  \draw[->] (0,0) -- (12.8,0);
  \node[anchor=west] at (12.8,0) {$t$\,(s)};
  \foreach \x in {0,...,12} {
    \draw (\x,0) -- (\x,-1.2);
    \node[below, font=\scriptsize\sffamily] at (\x,-1.2) {\x};
  }
  % arrivals
  \node[anchor=east, font=\scriptsize\sffamily] at (-0.2,-9) {\iflangja{到着}{arrival}};
  \foreach \x/\t in {0/T1, 1/T2, 2/T3} {
    \draw[->] (\x,-10.5) -- (\x,-5.5);
    \node[below, font=\scriptsize\sffamily, inner sep=1pt] at (\x,-10.5) {\t};
  }
\end{tikzpicture}
   (width ~100 mm)
\begin{tikzpicture}[x=7.5mm, y=1mm, font=\footnotesize\sffamily,
  >={Stealth[length=1.5mm]},
  t1/.style={draw, fill=ln-accent!25},
  t2/.style={draw, fill=ln-box},
  t3/.style={draw, fill=white},
  ttl/.style={anchor=south west, font=\footnotesize\sffamily\bfseries, inner sep=1pt}]
  % preemptive SJF
  \node[ttl] at (0,26.5) {\iflangja{(a) プリエンプティブ SJF}{(a) preemptive SJF}};
  \draw[t1] (0,20) rectangle (1,26);  \node at (0.5,23) {T1};
  \draw[t2] (1,20) rectangle (3,26);  \node at (2,23) {T2};
  \draw[t3] (3,20) rectangle (6,26);  \node at (4.5,23) {T3};
  \draw[t1] (6,20) rectangle (12,26); \node at (9,23) {T1};
  % priority
  \node[ttl] at (0,10.5) {\iflangja{(b) プリエンプティブ優先度スケジューリング}{(b) preemptive priority scheduling}};
  \draw[t1] (0,4) rectangle (2,10);   \node at (1,7) {T1};
  \draw[t3] (2,4) rectangle (5,10);   \node at (3.5,7) {T3};
  \draw[t1] (5,4) rectangle (10,10);  \node at (7.5,7) {T1};
  \draw[t2] (10,4) rectangle (12,10); \node at (11,7) {T2};
  % time axis
  \draw[->] (0,0) -- (12.8,0);
  \node[anchor=west] at (12.8,0) {$t$\,(s)};
  \foreach \x in {0,...,12} {
    \draw (\x,0) -- (\x,-1.2);
    \node[below, font=\scriptsize\sffamily] at (\x,-1.2) {\x};
  }
  % arrivals
  \node[anchor=east, font=\scriptsize\sffamily] at (-0.2,-9) {\iflangja{到着}{arrival}};
  \foreach \x/\t in {0/T1, 1/T2, 2/T3} {
    \draw[->] (\x,-10.5) -- (\x,-5.5);
    \node[below, font=\scriptsize\sffamily, inner sep=1pt] at (\x,-10.5) {\t};
  }
\end{tikzpicture}

  \caption{The tasks of \cref{ex:l07-preemptive} under (a) preemptive SJF
    and (b) preemptive priority scheduling.}
  \label{fig:l07-preemptive}
\end{figure}

\subsection{Starvation and aging}

Priority scheduling has a danger.  As long as tasks of higher priority keep
arriving, a task of low priority remains in the runqueue.

\begin{definition}[Starvation]
  A ready task that is indefinitely denied the CPU, because the scheduler
  always prefers other tasks, suffers \term{starvation}[スタベーション].
\end{definition}

The remedy is \term{priority aging}[エージング].  The priority used by the
scheduler is not the task's actual priority alone but a function of the
actual priority and of the time the task has spent in the runqueue.  The
longer a task waits, the higher its effective priority becomes, until it
exceeds that of the arriving tasks and the task runs.  Aging thus prevents
starvation.

\Needspace{7\baselineskip}
\section{Priority inversion}
\label{sec:l07-inversion}

Priorities interact with synchronization.\source{P3L1, priority inversion;
Silberschatz et al.\ ch.~5--6.}  Consider three tasks with priorities
$\text{T1} > \text{T2} > \text{T3}$, of which T1 and T3 lock the same mutex
\texttt{m} (\cref{fig:l07-inversion}a).  T3 arrives first and locks
\texttt{m} at time 1.  At time 2, T2 arrives and preempts T3; at time 3, T1
arrives and preempts T2.  At time 4, T1 tries to lock \texttt{m}, which is
held by T3, and blocks.  Of the two remaining tasks, T2 has the higher
priority, so T2 runs, and it runs to completion before T3 can continue.
Only then does T3 finish its critical section and unlock \texttt{m}, and T1
can proceed.  T2 completes at time 6, before T1, which completes at 9.

\begin{definition}[Priority inversion]
  In a \term{priority inversion}[優先度逆転] a task of high priority waits
  for a resource held by a task of lower priority, while tasks of
  intermediate priority, which do not need the resource, run in the
  meantime.  The high-priority task is delayed as if the priorities were
  inverted.
\end{definition}

\begin{figure}[htbp]
  \centering
  % Priority inversion and priority inheritance.  T1 (high) arrives at 3,
% T2 (medium) at 2, T3 (low) at 0; T1 and T3 share mutex m.
% Usage: % Priority inversion and priority inheritance.  T1 (high) arrives at 3,
% T2 (medium) at 2, T3 (low) at 0; T1 and T3 share mutex m.
% Usage: % Priority inversion and priority inheritance.  T1 (high) arrives at 3,
% T2 (medium) at 2, T3 (low) at 0; T1 and T3 share mutex m.
% Usage: \input{figures/l07-inversion}   (width ~112 mm)
\begin{tikzpicture}[x=9mm, y=1mm, font=\footnotesize\sffamily,
  run/.style={draw, fill=ln-box},
  cs/.style={draw, fill=ln-accent!40},
  boost/.style={draw=ln-caution, line width=1.2pt, fill=ln-accent!40},
  blk/.style={densely dashed, thick, ln-caution},
  lab/.style={anchor=east},
  ttl/.style={anchor=south west, font=\footnotesize\sffamily\bfseries, inner sep=1pt}]
  % ---------- (a) without priority inheritance
  \node[ttl] at (0,54) {\iflangja{(a) 優先度継承なし}{(a) without priority inheritance}};
  \node[lab] at (-0.1,50.5) {\iflangja{T1（高）}{T1 (high)}};
  \node[lab] at (-0.1,42.5) {\iflangja{T2（中）}{T2 (medium)}};
  \node[lab] at (-0.1,34.5) {\iflangja{T3（低）}{T3 (low)}};
  \draw[run] (3,48) rectangle (4,53);
  \draw[blk] (4,50.5) -- (7,50.5);
  \draw[cs]  (7,48) rectangle (8,53);
  \draw[run] (8,48) rectangle (9,53);
  \draw[run] (2,40) rectangle (3,45);
  \draw[run] (4,40) rectangle (6,45);
  \draw[run] (0,32) rectangle (1,37);
  \draw[cs]  (1,32) rectangle (2,37);
  \draw[cs]  (6,32) rectangle (7,37);
  \draw[run] (9,32) rectangle (10,37);
  \draw (0,29) -- (10.3,29);
  \foreach \x in {0,...,10} {
    \draw (\x,29) -- (\x,28);
    \node[below, font=\scriptsize\sffamily, inner sep=1pt] at (\x,28) {\x};
  }
  % ---------- (b) with priority inheritance
  \node[ttl] at (0,19) {\iflangja{(b) 優先度継承あり}{(b) with priority inheritance}};
  \node[lab] at (-0.1,15.5) {\iflangja{T1（高）}{T1 (high)}};
  \node[lab] at (-0.1,7.5)  {\iflangja{T2（中）}{T2 (medium)}};
  \node[lab] at (-0.1,-0.5) {\iflangja{T3（低）}{T3 (low)}};
  \draw[run] (3,13) rectangle (4,18);
  \draw[blk] (4,15.5) -- (5,15.5);
  \draw[cs]  (5,13) rectangle (6,18);
  \draw[run] (6,13) rectangle (7,18);
  \draw[run] (2,5) rectangle (3,10);
  \draw[run] (7,5) rectangle (9,10);
  \draw[run] (0,-3) rectangle (1,2);
  \draw[cs]  (1,-3) rectangle (2,2);
  \draw[boost] (4,-3) rectangle (5,2);
  \node[anchor=west, font=\scriptsize\sffamily, text=ln-caution] at (5.1,-0.5)
    {\iflangja{優先度を「高」に引き上げ}{priority raised to high}};
  \draw[run] (9,-3) rectangle (10,2);
  \draw (0,-6) -- (10.3,-6);
  \foreach \x in {0,...,10} {
    \draw (\x,-6) -- (\x,-7);
    \node[below, font=\scriptsize\sffamily, inner sep=1pt] at (\x,-7) {\x};
  }
  \node[anchor=west, font=\scriptsize\sffamily] at (10.3,-6) {$t$};
  % ---------- legend
  \begin{scope}[yshift=-17mm, font=\scriptsize\sffamily]
    \draw[run] (0,-1.5) rectangle (0.5,1.5);
    \node[anchor=west] at (0.55,0) {\iflangja{実行}{runs}};
    \draw[cs] (2,-1.5) rectangle (2.5,1.5);
    \node[anchor=west] at (2.55,0) {\iflangja{m を保持して実行}{runs holding m}};
    % The Japanese labels are wider: shift the last two entries so that the
    % "inherited priority" swatch does not cover the end of "blocked on m".
    \draw[blk] (\iflangja{5.0}{5.3},0) -- (\iflangja{5.6}{6},0);
    \node[anchor=west] at (\iflangja{5.65}{6.05},0) {\iflangja{m 待ちでブロック}{blocked on m}};
    \draw[boost] (\iflangja{8.45}{8.2},-1.5) rectangle (\iflangja{8.95}{8.7},1.5);
    \node[anchor=west] at (\iflangja{9.0}{8.75},0) {\iflangja{継承した優先度}{inherited priority}};
  \end{scope}
\end{tikzpicture}
   (width ~112 mm)
\begin{tikzpicture}[x=9mm, y=1mm, font=\footnotesize\sffamily,
  run/.style={draw, fill=ln-box},
  cs/.style={draw, fill=ln-accent!40},
  boost/.style={draw=ln-caution, line width=1.2pt, fill=ln-accent!40},
  blk/.style={densely dashed, thick, ln-caution},
  lab/.style={anchor=east},
  ttl/.style={anchor=south west, font=\footnotesize\sffamily\bfseries, inner sep=1pt}]
  % ---------- (a) without priority inheritance
  \node[ttl] at (0,54) {\iflangja{(a) 優先度継承なし}{(a) without priority inheritance}};
  \node[lab] at (-0.1,50.5) {\iflangja{T1（高）}{T1 (high)}};
  \node[lab] at (-0.1,42.5) {\iflangja{T2（中）}{T2 (medium)}};
  \node[lab] at (-0.1,34.5) {\iflangja{T3（低）}{T3 (low)}};
  \draw[run] (3,48) rectangle (4,53);
  \draw[blk] (4,50.5) -- (7,50.5);
  \draw[cs]  (7,48) rectangle (8,53);
  \draw[run] (8,48) rectangle (9,53);
  \draw[run] (2,40) rectangle (3,45);
  \draw[run] (4,40) rectangle (6,45);
  \draw[run] (0,32) rectangle (1,37);
  \draw[cs]  (1,32) rectangle (2,37);
  \draw[cs]  (6,32) rectangle (7,37);
  \draw[run] (9,32) rectangle (10,37);
  \draw (0,29) -- (10.3,29);
  \foreach \x in {0,...,10} {
    \draw (\x,29) -- (\x,28);
    \node[below, font=\scriptsize\sffamily, inner sep=1pt] at (\x,28) {\x};
  }
  % ---------- (b) with priority inheritance
  \node[ttl] at (0,19) {\iflangja{(b) 優先度継承あり}{(b) with priority inheritance}};
  \node[lab] at (-0.1,15.5) {\iflangja{T1（高）}{T1 (high)}};
  \node[lab] at (-0.1,7.5)  {\iflangja{T2（中）}{T2 (medium)}};
  \node[lab] at (-0.1,-0.5) {\iflangja{T3（低）}{T3 (low)}};
  \draw[run] (3,13) rectangle (4,18);
  \draw[blk] (4,15.5) -- (5,15.5);
  \draw[cs]  (5,13) rectangle (6,18);
  \draw[run] (6,13) rectangle (7,18);
  \draw[run] (2,5) rectangle (3,10);
  \draw[run] (7,5) rectangle (9,10);
  \draw[run] (0,-3) rectangle (1,2);
  \draw[cs]  (1,-3) rectangle (2,2);
  \draw[boost] (4,-3) rectangle (5,2);
  \node[anchor=west, font=\scriptsize\sffamily, text=ln-caution] at (5.1,-0.5)
    {\iflangja{優先度を「高」に引き上げ}{priority raised to high}};
  \draw[run] (9,-3) rectangle (10,2);
  \draw (0,-6) -- (10.3,-6);
  \foreach \x in {0,...,10} {
    \draw (\x,-6) -- (\x,-7);
    \node[below, font=\scriptsize\sffamily, inner sep=1pt] at (\x,-7) {\x};
  }
  \node[anchor=west, font=\scriptsize\sffamily] at (10.3,-6) {$t$};
  % ---------- legend
  \begin{scope}[yshift=-17mm, font=\scriptsize\sffamily]
    \draw[run] (0,-1.5) rectangle (0.5,1.5);
    \node[anchor=west] at (0.55,0) {\iflangja{実行}{runs}};
    \draw[cs] (2,-1.5) rectangle (2.5,1.5);
    \node[anchor=west] at (2.55,0) {\iflangja{m を保持して実行}{runs holding m}};
    % The Japanese labels are wider: shift the last two entries so that the
    % "inherited priority" swatch does not cover the end of "blocked on m".
    \draw[blk] (\iflangja{5.0}{5.3},0) -- (\iflangja{5.6}{6},0);
    \node[anchor=west] at (\iflangja{5.65}{6.05},0) {\iflangja{m 待ちでブロック}{blocked on m}};
    \draw[boost] (\iflangja{8.45}{8.2},-1.5) rectangle (\iflangja{8.95}{8.7},1.5);
    \node[anchor=west] at (\iflangja{9.0}{8.75},0) {\iflangja{継承した優先度}{inherited priority}};
  \end{scope}
\end{tikzpicture}
   (width ~112 mm)
\begin{tikzpicture}[x=9mm, y=1mm, font=\footnotesize\sffamily,
  run/.style={draw, fill=ln-box},
  cs/.style={draw, fill=ln-accent!40},
  boost/.style={draw=ln-caution, line width=1.2pt, fill=ln-accent!40},
  blk/.style={densely dashed, thick, ln-caution},
  lab/.style={anchor=east},
  ttl/.style={anchor=south west, font=\footnotesize\sffamily\bfseries, inner sep=1pt}]
  % ---------- (a) without priority inheritance
  \node[ttl] at (0,54) {\iflangja{(a) 優先度継承なし}{(a) without priority inheritance}};
  \node[lab] at (-0.1,50.5) {\iflangja{T1（高）}{T1 (high)}};
  \node[lab] at (-0.1,42.5) {\iflangja{T2（中）}{T2 (medium)}};
  \node[lab] at (-0.1,34.5) {\iflangja{T3（低）}{T3 (low)}};
  \draw[run] (3,48) rectangle (4,53);
  \draw[blk] (4,50.5) -- (7,50.5);
  \draw[cs]  (7,48) rectangle (8,53);
  \draw[run] (8,48) rectangle (9,53);
  \draw[run] (2,40) rectangle (3,45);
  \draw[run] (4,40) rectangle (6,45);
  \draw[run] (0,32) rectangle (1,37);
  \draw[cs]  (1,32) rectangle (2,37);
  \draw[cs]  (6,32) rectangle (7,37);
  \draw[run] (9,32) rectangle (10,37);
  \draw (0,29) -- (10.3,29);
  \foreach \x in {0,...,10} {
    \draw (\x,29) -- (\x,28);
    \node[below, font=\scriptsize\sffamily, inner sep=1pt] at (\x,28) {\x};
  }
  % ---------- (b) with priority inheritance
  \node[ttl] at (0,19) {\iflangja{(b) 優先度継承あり}{(b) with priority inheritance}};
  \node[lab] at (-0.1,15.5) {\iflangja{T1（高）}{T1 (high)}};
  \node[lab] at (-0.1,7.5)  {\iflangja{T2（中）}{T2 (medium)}};
  \node[lab] at (-0.1,-0.5) {\iflangja{T3（低）}{T3 (low)}};
  \draw[run] (3,13) rectangle (4,18);
  \draw[blk] (4,15.5) -- (5,15.5);
  \draw[cs]  (5,13) rectangle (6,18);
  \draw[run] (6,13) rectangle (7,18);
  \draw[run] (2,5) rectangle (3,10);
  \draw[run] (7,5) rectangle (9,10);
  \draw[run] (0,-3) rectangle (1,2);
  \draw[cs]  (1,-3) rectangle (2,2);
  \draw[boost] (4,-3) rectangle (5,2);
  \node[anchor=west, font=\scriptsize\sffamily, text=ln-caution] at (5.1,-0.5)
    {\iflangja{優先度を「高」に引き上げ}{priority raised to high}};
  \draw[run] (9,-3) rectangle (10,2);
  \draw (0,-6) -- (10.3,-6);
  \foreach \x in {0,...,10} {
    \draw (\x,-6) -- (\x,-7);
    \node[below, font=\scriptsize\sffamily, inner sep=1pt] at (\x,-7) {\x};
  }
  \node[anchor=west, font=\scriptsize\sffamily] at (10.3,-6) {$t$};
  % ---------- legend
  \begin{scope}[yshift=-17mm, font=\scriptsize\sffamily]
    \draw[run] (0,-1.5) rectangle (0.5,1.5);
    \node[anchor=west] at (0.55,0) {\iflangja{実行}{runs}};
    \draw[cs] (2,-1.5) rectangle (2.5,1.5);
    \node[anchor=west] at (2.55,0) {\iflangja{m を保持して実行}{runs holding m}};
    % The Japanese labels are wider: shift the last two entries so that the
    % "inherited priority" swatch does not cover the end of "blocked on m".
    \draw[blk] (\iflangja{5.0}{5.3},0) -- (\iflangja{5.6}{6},0);
    \node[anchor=west] at (\iflangja{5.65}{6.05},0) {\iflangja{m 待ちでブロック}{blocked on m}};
    \draw[boost] (\iflangja{8.45}{8.2},-1.5) rectangle (\iflangja{8.95}{8.7},1.5);
    \node[anchor=west] at (\iflangja{9.0}{8.75},0) {\iflangja{継承した優先度}{inherited priority}};
  \end{scope}
\end{tikzpicture}

  \caption{T1 and T3 share mutex \texttt{m}.  (a) While T1 is blocked on
    \texttt{m}, the medium-priority T2 runs to completion ahead of it.
    (b) With priority inheritance, T3 runs at the priority of T1 until it
    unlocks \texttt{m}, and T1 completes before T2.}
  \label{fig:l07-inversion}
\end{figure}

The solution is to raise the priority of the owner of the mutex temporarily
(\cref{fig:l07-inversion}b).  When T1 blocks on \texttt{m}, T3 is given the
priority of T1, so that T3, not T2, runs and finishes its critical section.
When T3 unlocks \texttt{m}, its priority is lowered back to its own, T1
acquires the mutex and runs, and T1 completes at time 7, ahead of T2.

\begin{definition}[Priority inheritance]
  With \term{priority inheritance}[優先度継承] the owner of a lock
  temporarily runs at the priority of the highest-priority task that is
  blocked on the lock, and returns to its own priority when it releases the
  lock.
\end{definition}

\Needspace{9\baselineskip}
To boost the owner, the lock operation must know who the owner is; this is
another reason for recording the owner in the mutex data structure
(Lesson~5).\sidenote{In July 1997 the Mars Pathfinder lander reset itself
repeatedly: a low-priority meteorological task held the mutex of the shared
information bus, the high-priority bus task blocked on it, and
medium-priority communication ran instead until a watchdog timer reset the
spacecraft.  JPL reproduced the inversion on the ground and uploaded a patch
that enabled priority inheritance on that mutex.}  POSIX threads (Lesson~4)
provide priority inheritance as a protocol attribute of a
mutex:
\begin{lstlisting}[language=C, numbers=none]
pthread_mutexattr_t attr;
pthread_mutex_t m;

pthread_mutexattr_init(&attr);
pthread_mutexattr_setprotocol(&attr, PTHREAD_PRIO_INHERIT);
pthread_mutex_init(&m, &attr);
\end{lstlisting}

\begin{keypoint}
  Preemption lets a scheduler serve short or important tasks first.
  Priorities bring the risks of starvation, avoided by aging, and of
  priority inversion, avoided by letting a lock owner inherit the priority
  of the tasks blocked on its lock.\caution{POSIX also defines
  \texttt{PTHREAD\_PRIO\_PROTECT}, the priority-ceiling protocol: the owner
  runs at a ceiling priority declared for the mutex, whether or not anyone
  is blocked.  The lecture treats only inheritance.}
\end{keypoint}

\Needspace{7\baselineskip}
\section{Round robin scheduling and timesharing}
\label{sec:l07-rr}

A scheduler for tasks whose execution times are unknown can simply let them
take turns.\source{P3L1, round robin scheduling, timesharing and
timeslices; OSTEP ch.~7.}  Such a scheduler needs neither estimates of
execution times nor, in its basic form, priorities.

\begin{definition}[Round robin]
  A scheduler that uses
  \term{round robin scheduling}[ラウンドロビンスケジューリング] (RR) takes
  the task at the head of the runqueue, like FCFS.  Unlike FCFS, a task
  need not run to completion: it may yield the CPU, for example to wait
  for I/O, and when it is ready again it is placed at the tail of the
  runqueue.
\end{definition}

Round robin can be combined with priorities: tasks of equal priority take
turns, and a task of higher priority that becomes ready preempts the
running one.  In its most common form, round robin interleaves the tasks
even when they do not yield, by limiting how long each task may run.

The timeslice of Lesson~2\index{timeslice} is the maximum amount of
uninterrupted time given to a task; it is also called the time
quantum.\margin{Linux still gives its real-time round-robin class a fixed
quantum: \SI{100}{\milli\second} by default, set with the sysctl
\texttt{sched\_rr\_timeslice\_ms}.}  A
task may run for less than its timeslice.  It may block, waiting for I/O or
on a synchronization construct, and be placed on the corresponding queue;
or a task of higher priority may become ready and preempt it.  When the
timeslice expires, the running task is preempted and returns to the tail
of the runqueue.

\begin{definition}[Timesharing]
  The interleaving of tasks by means of timeslices is called
  \term{timesharing}[タイムシェアリング] of the CPU: every task runs for at
  most one timeslice at a time, so that a CPU-intensive task is preempted
  when its timeslice expires and cannot hold the CPU until it completes.
\end{definition}

\begin{example}
  The tasks of \cref{ex:l07-rtc} run under round robin with a timeslice of
  \SI{1}{\second}, ignoring the cost of context switches
  (\cref{fig:l07-gantt}).  The order is T1, T2, T3, T1, T2, T1: T3
  completes at \SI{3}{\second}, T2 at \SI{5}{\second} and T1 at
  \SI{8}{\second}, after running alone for its last three timeslices.  The
  metrics are compared in \cref{tab:l07-rr}.  Without knowing any execution
  time, round robin lowers the average completion time from
  \SI{6.67}{\second} to \SI{5.33}{\second}, halfway from FCFS to the
  \SI{4.00}{\second} of SJF, and since a task waits for at most one
  timeslice of each task ahead of it before it first runs, its average wait
  time is even shorter than that of SJF.
\end{example}

\begin{table}[htbp]
  \centering
  \caption{Metrics of the three schedules of \cref{fig:l07-gantt}.}
  \label{tab:l07-rr}
  \small
  \begin{tabular}{@{}lccc@{}}
    \toprule
    & FCFS & SJF & RR (\SI{1}{\second}) \\
    \midrule
    Throughput (tasks/s)    & 0.375 & 0.375 & 0.375 \\
    Average completion time (s) & 6.67 & 4.00 & 5.33 \\
    Average wait time (s)       & 4.00 & 1.33 & 1.00 \\
    \bottomrule
  \end{tabular}
\end{table}

Timesharing has several benefits: short tasks complete sooner; the system
is more responsive, since every ready task runs within a bounded time; and
tasks that initiate lengthy I/O operations reach them sooner, so that the
operations can proceed while other tasks compute.  It also has a cost that
the example ignores.  At the end of every timeslice the timer interrupts
the running task, the scheduler runs, and a context switch takes place.  To
keep this overhead small, the timeslice must be much longer than the time
taken by the scheduler and the context switch, as \cref{eq:l02-useful}
showed for the time of the scheduler alone.

\Needspace{7\baselineskip}
\section{Choosing the timeslice length}
\label{sec:l07-timeslice}

How long a timeslice should be depends on how the tasks use
it.\source{P3L1, how long should a timeslice be, CPU bound timeslice
length, I/O bound timeslice length, summarizing timeslice length; OSTEP
ch.~7.}  The timeslice must balance the benefits of timesharing against its
overhead, and the balance differs between tasks that mostly compute and
tasks that mostly wait for I/O.

\subsection{CPU-bound tasks}

A task that spends most of its time computing and rarely blocks is a
\term{CPU-bound task}[CPUバウンドタスク].

\begin{example}
  \label{ex:l07-cpu-bound}
  Two CPU-bound tasks, T1 and T2, each need \SI{6}{\second} of CPU time and
  arrive at time 0; a context switch takes \SI{0.2}{\second}.\caution{The
  \SI{0.2}{\second} is an illustration, not a measurement: Lesson~2 puts the
  direct cost of a switch at a few microseconds, so a real system loses far
  less per timeslice than this example suggests.}  With a
  timeslice of \SI{1}{\second} the tasks alternate twelve times, with eleven
  context switches between the timeslices.  T2 starts at \SI{1.2}{\second}.
  The last timeslice of T1 starts after ten timeslices and ten context
  switches, at \SI{12}{\second}, so T1 completes at \SI{13}{\second} and T2,
  after one more context switch, at \SI{14.2}{\second}.  With a timeslice of \SI{3}{\second}
  the tasks alternate four times, with three context switches: T2 starts
  at \SI{3.2}{\second}, T1 completes at \SI{9.4}{\second} and T2 at
  \SI{12.6}{\second}.  With an infinite timeslice T1 runs to completion at
  \SI{6}{\second} and T2 runs from 6.2 to \SI{12.2}{\second}.
  \Cref{tab:l07-cpu-bound} lists the resulting metrics.
\end{example}

\begin{table}[htbp]
  \centering
  \caption{Two CPU-bound tasks of \SI{6}{\second} with a context switch
    of \SI{0.2}{\second} (\cref{ex:l07-cpu-bound}).}
  \label{tab:l07-cpu-bound}
  \small
  \begin{tabular}{@{}lccc@{}}
    \toprule
    Timeslice & Throughput (tasks/s) & Avg.\ wait (s) & Avg.\ completion (s) \\
    \midrule
    \SI{1}{\second}  & $2/14.2 \approx 0.141$ & $(0+1.2)/2 = 0.6$ & $(13+14.2)/2 = 13.6$ \\
    \SI{3}{\second}  & $2/12.6 \approx 0.159$ & $(0+3.2)/2 = 1.6$ & $(9.4+12.6)/2 = 11.0$ \\
    $\infty$         & $2/12.2 \approx 0.164$ & $(0+6.2)/2 = 3.1$ & $(6+12.2)/2 = 9.1$ \\
    \bottomrule
  \end{tabular}
\end{table}

For CPU-bound tasks, a longer timeslice improves throughput and average
completion time, because fewer context switches are spent.  Only the wait
time grows, and for CPU-bound tasks, such as batch computations, it matters
little: their users see the time at which a computation completes, not the
time at which it starts.  In the limit, an infinite timeslice, that is,
run to completion, is best for them.

\subsection{I/O-bound tasks}

A task that frequently blocks for I/O and computes only briefly between its
I/O operations is an \term{I/O-bound task}[I/Oバウンドタスク].  If all tasks
are I/O-bound, the length of the timeslice hardly matters, as long as it
exceeds their short bursts of computation: each task then blocks for I/O
before its timeslice expires and leaves the CPU of its own accord.  The length matters when I/O-bound tasks share the CPU with
CPU-bound ones.

\begin{figure}[htbp]
  \centering
  % A CPU-bound task T1 and an I/O-bound task T2 (5 ms of computation, then
% an I/O operation of 10 ms) under round robin with a context switch of
% 2 ms and timeslices of 5 ms and 30 ms.
% Usage: % A CPU-bound task T1 and an I/O-bound task T2 (5 ms of computation, then
% an I/O operation of 10 ms) under round robin with a context switch of
% 2 ms and timeslices of 5 ms and 30 ms.
% Usage: % A CPU-bound task T1 and an I/O-bound task T2 (5 ms of computation, then
% an I/O operation of 10 ms) under round robin with a context switch of
% 2 ms and timeslices of 5 ms and 30 ms.
% Usage: \input{figures/l07-io-timeslice}   (width ~104 mm)
\begin{tikzpicture}[x=1.1mm, y=1mm, font=\scriptsize\sffamily,
  >={Stealth[length=1.4mm]},
  t1/.style={draw, fill=ln-accent!25},
  t2/.style={draw, fill=white},
  cx/.style={draw, fill=ln-muted},
  io/.style={draw, fill=ln-box},
  lab/.style={anchor=east, font=\footnotesize\sffamily},
  ttl/.style={anchor=south west, font=\footnotesize\sffamily\bfseries, inner sep=1pt}]
  % ---------- (a) timeslice 5 ms
  \node[ttl] at (0,45) {\iflangja{(a) タイムスライス 5\,ms}{(a) timeslice 5\,ms}};
  \node[lab] at (-1,41) {CPU};
  \node[lab] at (-1,31) {I/O};
  \foreach \a in {0,19,38,57} {
    \draw[t2] (\a,38) rectangle (\a+5,44);
    \node at (\a+2.5,41) {T2};
    \draw[cx] (\a+5,38) rectangle (\a+7,44);
    \draw[io] (\a+5,28) rectangle (\a+15,34);
    \node at (\a+10,31) {T2};
  }
  \foreach \a in {7,26,45,64} {
    \draw[t1] (\a,38) rectangle (\a+10,44);
    \draw[densely dotted] (\a+5,38) -- (\a+5,44);
    \node at (\a+2.5,41) {T1};
    \node at (\a+7.5,41) {T1};
    \draw[cx] (\a+10,38) rectangle (\a+12,44);
  }
  \draw[t2] (76,38) -- (80,38) (76,44) -- (80,44) (76,38) -- (76,44);
  \node at (78.2,41) {T2};
  \draw (0,25) -- (82,25);
  \foreach \x in {0,10,...,80} {
    \draw (\x,25) -- (\x,24);
    \node[below, inner sep=1pt] at (\x,24) {\x};
  }
  % ---------- (b) timeslice 30 ms
  \node[ttl] at (0,13) {\iflangja{(b) タイムスライス 30\,ms}{(b) timeslice 30\,ms}};
  \node[lab] at (-1,9) {CPU};
  \node[lab] at (-1,-7) {I/O};
  \foreach \a in {0,39} {
    \draw[t2] (\a,6) rectangle (\a+5,12);
    \node at (\a+2.5,9) {T2};
    \draw[cx] (\a+5,6) rectangle (\a+7,12);
    \draw[io] (\a+5,-10) rectangle (\a+15,-4);
    \node at (\a+10,-7) {T2};
    \draw[t1] (\a+7,6) rectangle (\a+37,12);
    \node at (\a+22,9) {T1};
    \draw[cx] (\a+37,6) rectangle (\a+39,12);
  }
  \draw[t2] (78,6) -- (80,6) (78,12) -- (80,12) (78,6) -- (78,12);
  % T2 is ready but waits for the timeslice of T1 to end
  \draw[<->, ln-caution] (15,2.5) -- (39,2.5);
  \draw[ln-caution, densely dotted] (15,-4) -- (15,2.5) (39,2.5) -- (39,6);
  \node[anchor=north west, text=ln-caution, inner sep=1pt] at (16,1.8)
    {\iflangja{T2 は実行可能だが待つ（24\,ms）}{T2 ready, waiting (24\,ms)}};
  \draw (0,-13) -- (82,-13);
  \foreach \x in {0,10,...,80} {
    \draw (\x,-13) -- (\x,-14);
    \node[below, inner sep=1pt] at (\x,-14) {\x};
  }
  \node[anchor=west] at (82,-13) {ms};
  % ---------- legend
  \begin{scope}[yshift=-24mm]
    \draw[t1] (0,-1.5) rectangle (4,1.5);
    \node[anchor=west] at (4.5,0) {\iflangja{CPU バウンドの T1}{CPU-bound T1}};
    \draw[t2] (28,-1.5) rectangle (32,1.5);
    \node[anchor=west] at (32.5,0) {\iflangja{I/O バウンドの T2}{I/O-bound T2}};
    \draw[cx] (55,-1.5) rectangle (57,1.5);
    \node[anchor=west] at (57.5,0) {\iflangja{コンテキストスイッチ}{context switch}};
  \end{scope}
\end{tikzpicture}
   (width ~104 mm)
\begin{tikzpicture}[x=1.1mm, y=1mm, font=\scriptsize\sffamily,
  >={Stealth[length=1.4mm]},
  t1/.style={draw, fill=ln-accent!25},
  t2/.style={draw, fill=white},
  cx/.style={draw, fill=ln-muted},
  io/.style={draw, fill=ln-box},
  lab/.style={anchor=east, font=\footnotesize\sffamily},
  ttl/.style={anchor=south west, font=\footnotesize\sffamily\bfseries, inner sep=1pt}]
  % ---------- (a) timeslice 5 ms
  \node[ttl] at (0,45) {\iflangja{(a) タイムスライス 5\,ms}{(a) timeslice 5\,ms}};
  \node[lab] at (-1,41) {CPU};
  \node[lab] at (-1,31) {I/O};
  \foreach \a in {0,19,38,57} {
    \draw[t2] (\a,38) rectangle (\a+5,44);
    \node at (\a+2.5,41) {T2};
    \draw[cx] (\a+5,38) rectangle (\a+7,44);
    \draw[io] (\a+5,28) rectangle (\a+15,34);
    \node at (\a+10,31) {T2};
  }
  \foreach \a in {7,26,45,64} {
    \draw[t1] (\a,38) rectangle (\a+10,44);
    \draw[densely dotted] (\a+5,38) -- (\a+5,44);
    \node at (\a+2.5,41) {T1};
    \node at (\a+7.5,41) {T1};
    \draw[cx] (\a+10,38) rectangle (\a+12,44);
  }
  \draw[t2] (76,38) -- (80,38) (76,44) -- (80,44) (76,38) -- (76,44);
  \node at (78.2,41) {T2};
  \draw (0,25) -- (82,25);
  \foreach \x in {0,10,...,80} {
    \draw (\x,25) -- (\x,24);
    \node[below, inner sep=1pt] at (\x,24) {\x};
  }
  % ---------- (b) timeslice 30 ms
  \node[ttl] at (0,13) {\iflangja{(b) タイムスライス 30\,ms}{(b) timeslice 30\,ms}};
  \node[lab] at (-1,9) {CPU};
  \node[lab] at (-1,-7) {I/O};
  \foreach \a in {0,39} {
    \draw[t2] (\a,6) rectangle (\a+5,12);
    \node at (\a+2.5,9) {T2};
    \draw[cx] (\a+5,6) rectangle (\a+7,12);
    \draw[io] (\a+5,-10) rectangle (\a+15,-4);
    \node at (\a+10,-7) {T2};
    \draw[t1] (\a+7,6) rectangle (\a+37,12);
    \node at (\a+22,9) {T1};
    \draw[cx] (\a+37,6) rectangle (\a+39,12);
  }
  \draw[t2] (78,6) -- (80,6) (78,12) -- (80,12) (78,6) -- (78,12);
  % T2 is ready but waits for the timeslice of T1 to end
  \draw[<->, ln-caution] (15,2.5) -- (39,2.5);
  \draw[ln-caution, densely dotted] (15,-4) -- (15,2.5) (39,2.5) -- (39,6);
  \node[anchor=north west, text=ln-caution, inner sep=1pt] at (16,1.8)
    {\iflangja{T2 は実行可能だが待つ（24\,ms）}{T2 ready, waiting (24\,ms)}};
  \draw (0,-13) -- (82,-13);
  \foreach \x in {0,10,...,80} {
    \draw (\x,-13) -- (\x,-14);
    \node[below, inner sep=1pt] at (\x,-14) {\x};
  }
  \node[anchor=west] at (82,-13) {ms};
  % ---------- legend
  \begin{scope}[yshift=-24mm]
    \draw[t1] (0,-1.5) rectangle (4,1.5);
    \node[anchor=west] at (4.5,0) {\iflangja{CPU バウンドの T1}{CPU-bound T1}};
    \draw[t2] (28,-1.5) rectangle (32,1.5);
    \node[anchor=west] at (32.5,0) {\iflangja{I/O バウンドの T2}{I/O-bound T2}};
    \draw[cx] (55,-1.5) rectangle (57,1.5);
    \node[anchor=west] at (57.5,0) {\iflangja{コンテキストスイッチ}{context switch}};
  \end{scope}
\end{tikzpicture}
   (width ~104 mm)
\begin{tikzpicture}[x=1.1mm, y=1mm, font=\scriptsize\sffamily,
  >={Stealth[length=1.4mm]},
  t1/.style={draw, fill=ln-accent!25},
  t2/.style={draw, fill=white},
  cx/.style={draw, fill=ln-muted},
  io/.style={draw, fill=ln-box},
  lab/.style={anchor=east, font=\footnotesize\sffamily},
  ttl/.style={anchor=south west, font=\footnotesize\sffamily\bfseries, inner sep=1pt}]
  % ---------- (a) timeslice 5 ms
  \node[ttl] at (0,45) {\iflangja{(a) タイムスライス 5\,ms}{(a) timeslice 5\,ms}};
  \node[lab] at (-1,41) {CPU};
  \node[lab] at (-1,31) {I/O};
  \foreach \a in {0,19,38,57} {
    \draw[t2] (\a,38) rectangle (\a+5,44);
    \node at (\a+2.5,41) {T2};
    \draw[cx] (\a+5,38) rectangle (\a+7,44);
    \draw[io] (\a+5,28) rectangle (\a+15,34);
    \node at (\a+10,31) {T2};
  }
  \foreach \a in {7,26,45,64} {
    \draw[t1] (\a,38) rectangle (\a+10,44);
    \draw[densely dotted] (\a+5,38) -- (\a+5,44);
    \node at (\a+2.5,41) {T1};
    \node at (\a+7.5,41) {T1};
    \draw[cx] (\a+10,38) rectangle (\a+12,44);
  }
  \draw[t2] (76,38) -- (80,38) (76,44) -- (80,44) (76,38) -- (76,44);
  \node at (78.2,41) {T2};
  \draw (0,25) -- (82,25);
  \foreach \x in {0,10,...,80} {
    \draw (\x,25) -- (\x,24);
    \node[below, inner sep=1pt] at (\x,24) {\x};
  }
  % ---------- (b) timeslice 30 ms
  \node[ttl] at (0,13) {\iflangja{(b) タイムスライス 30\,ms}{(b) timeslice 30\,ms}};
  \node[lab] at (-1,9) {CPU};
  \node[lab] at (-1,-7) {I/O};
  \foreach \a in {0,39} {
    \draw[t2] (\a,6) rectangle (\a+5,12);
    \node at (\a+2.5,9) {T2};
    \draw[cx] (\a+5,6) rectangle (\a+7,12);
    \draw[io] (\a+5,-10) rectangle (\a+15,-4);
    \node at (\a+10,-7) {T2};
    \draw[t1] (\a+7,6) rectangle (\a+37,12);
    \node at (\a+22,9) {T1};
    \draw[cx] (\a+37,6) rectangle (\a+39,12);
  }
  \draw[t2] (78,6) -- (80,6) (78,12) -- (80,12) (78,6) -- (78,12);
  % T2 is ready but waits for the timeslice of T1 to end
  \draw[<->, ln-caution] (15,2.5) -- (39,2.5);
  \draw[ln-caution, densely dotted] (15,-4) -- (15,2.5) (39,2.5) -- (39,6);
  \node[anchor=north west, text=ln-caution, inner sep=1pt] at (16,1.8)
    {\iflangja{T2 は実行可能だが待つ（24\,ms）}{T2 ready, waiting (24\,ms)}};
  \draw (0,-13) -- (82,-13);
  \foreach \x in {0,10,...,80} {
    \draw (\x,-13) -- (\x,-14);
    \node[below, inner sep=1pt] at (\x,-14) {\x};
  }
  \node[anchor=west] at (82,-13) {ms};
  % ---------- legend
  \begin{scope}[yshift=-24mm]
    \draw[t1] (0,-1.5) rectangle (4,1.5);
    \node[anchor=west] at (4.5,0) {\iflangja{CPU バウンドの T1}{CPU-bound T1}};
    \draw[t2] (28,-1.5) rectangle (32,1.5);
    \node[anchor=west] at (32.5,0) {\iflangja{I/O バウンドの T2}{I/O-bound T2}};
    \draw[cx] (55,-1.5) rectangle (57,1.5);
    \node[anchor=west] at (57.5,0) {\iflangja{コンテキストスイッチ}{context switch}};
  \end{scope}
\end{tikzpicture}

  \caption{A CPU-bound task T1 and an I/O-bound task T2 under round robin
    (\cref{ex:l07-io-bound}).  When a timeslice of T1 expires while T2 is
    still waiting for its I/O, T1 simply continues.}
  \label{fig:l07-io-timeslice}
\end{figure}

\begin{example}
  \label{ex:l07-io-bound}
  A CPU-bound task T1 shares the CPU with an I/O-bound task T2, which
  computes for \SI{5}{\milli\second} and then issues an I/O operation that
  takes \SI{10}{\milli\second}; a context switch takes
  \SI{2}{\milli\second} (\cref{fig:l07-io-timeslice}).  With a timeslice of
  \SI{5}{\milli\second}, the I/O operation of T2 completes during the second
  timeslice of T1, and T2 waits \SI{4}{\milli\second} before it runs again.
  The pattern repeats every \SI{19}{\milli\second}: T2 issues an I/O
  operation every \SI{19}{\milli\second}, and the device is busy for
  $10/19 \approx \SI{53}{\percent}$ of the time.  With a timeslice of
  \SI{30}{\milli\second}, T2 waits \SI{24}{\milli\second} after each I/O
  operation, the pattern repeats every \SI{39}{\milli\second}, and the
  device is busy only $10/39 \approx \SI{26}{\percent}$ of the time.  T1, in
  turn, receives \SI{10}{\milli\second} of every \SI{19}{\milli\second}
  ($\approx\SI{53}{\percent}$) with the short timeslice but
  \SI{30}{\milli\second} of every \SI{39}{\milli\second}
  ($\approx\SI{77}{\percent}$) with the long one.
\end{example}

\subsection{Summary}

The two kinds of tasks prefer opposite timeslice lengths.
\begin{description}
  \item[CPU-bound tasks prefer longer timeslices.]  Fewer context switches
    limit the overhead and keep CPU utilization and throughput high.
  \item[I/O-bound tasks prefer shorter timeslices.]  An I/O-bound task runs
    soon after its I/O completes and issues its next I/O operation earlier,
    which keeps the devices busy and gives better user-perceived
    performance, since interactive tasks are typically
    I/O-bound.\margin{The \SI{10}{\milli\second} of the example is a disk
    figure: a \num{7200}~rpm drive averages about \SI{4}{\milli\second} of
    rotation and a seek of similar length, whereas an NVMe SSD answers a read
    in tens of microseconds.}
\end{description}

\subsection{CPU utilization}

The fraction of CPU time spent on useful work, \cref{eq:l02-useful},
extends to a mix of tasks with different behaviour.\tip{For one switch of
cost $c$ per timeslice $q$ the fraction is $q/(q+c)$: to keep the loss below
\SI{1}{\percent}, give a timeslice at least a hundred times the cost of a
context switch.}  Choose a recurring
interval of the schedule in which every task has had its opportunity to
run.  Within that interval, let $t_{\mathrm{run}}$ be the total time the CPU
spends running tasks and $t_{\mathrm{ctx}}$ the total time it spends on
context switches.  Then
\begin{equation}
  \text{CPU utilization} \;=\;
  \frac{t_{\mathrm{run}}}{t_{\mathrm{run}} + t_{\mathrm{ctx}}} .
  \label{eq:l07-utilization}
\end{equation}

\begin{example}
  Three I/O-bound tasks each compute for \SI{2}{\milli\second} and then
  issue an I/O operation of \SI{5}{\milli\second}; two CPU-bound tasks
  never block; a context switch takes \SI{0.5}{\milli\second}; all tasks run
  for a long time.  With a timeslice of \SI{2}{\milli\second}, every task,
  CPU-bound or not, leaves the CPU after \SI{2}{\milli\second}, and every
  run is followed by a context switch, so the utilization is
  $2/(2+0.5) = \SI{80}{\percent}$.  With a timeslice of
  \SI{20}{\milli\second}, one round consists of $3\cdot 2 + 2\cdot 20 =
  \SI{46}{\milli\second}$ of computation and five context switches; the
  utilization is $46/(46+2.5) \approx \SI{94.8}{\percent}$.  Both rounds,
  of 12.5 and \SI{48.5}{\milli\second}, are longer than an I/O operation,
  so the I/O-bound tasks are always ready when their turn
  comes.\caution{\Cref{eq:l07-utilization} ignores idle time: a stretch in
  which every task waits for I/O adds to neither $t_{\mathrm{run}}$ nor
  $t_{\mathrm{ctx}}$, so the ratio can read \SI{100}{\percent} while the CPU
  often has nothing to run.}
\end{example}

\begin{keypoint}
  Timesharing makes a system responsive without knowledge of execution
  times, at the price of an interrupt, a scheduler run and a context
  switch per timeslice.  CPU-bound tasks favour long timeslices, I/O-bound
  tasks short ones.
\end{keypoint}

\Needspace{7\baselineskip}
\section{Multi-level feedback queues}
\label{sec:l07-mlfq}

To give I/O-bound and CPU-bound tasks different timeslices, the scheduler
can keep all tasks in one runqueue and check the type of each task when it
is selected, or it can keep the two kinds of tasks in different
structures.\source{P3L1, runqueue data structure; OSTEP ch.~8; Corbató et
al.\ \cite{corbato1962}.}  A runqueue of several queues takes the second
approach (\cref{fig:l07-mlfq}):
\begin{itemize}
  \item the first queue holds the most I/O-intensive tasks, has the highest
    priority, and gives its tasks a timeslice of \SI{8}{\milli\second};
  \item the second queue holds tasks that mix I/O and computation, with a
    timeslice of \SI{16}{\milli\second};
  \item the last queue holds the CPU-intensive tasks, has the lowest
    priority, and gives them an infinite timeslice, that is, it schedules
    them FCFS.
\end{itemize}
The I/O-bound tasks receive the benefits of timesharing, and the CPU-bound
tasks avoid its overhead.

\begin{figure}[htbp]
  \centering
  % Multi-level feedback queue with three levels.
% Usage: % Multi-level feedback queue with three levels.
% Usage: % Multi-level feedback queue with three levels.
% Usage: \input{figures/l07-mlfq}   (width ~104 mm)
\begin{tikzpicture}[x=1mm, y=1mm, font=\footnotesize\sffamily,
  >={Stealth[length=1.6mm]},
  desc/.style={anchor=east, align=right, font=\scriptsize\sffamily},
  q/.style={draw, fill=ln-box},
  tk/.style={draw, circle, fill=white, minimum size=3.2mm, inner sep=0pt}]
  \foreach \y in {28,14,0} {
    \draw[q] (0,\y-3) rectangle (36,\y+3);
    \foreach \i in {1,...,4} \draw (36-5*\i,\y-3) -- ++(0,6);
    \draw[->] (36,\y) -- (44,\y);
  }
  \node[tk] at (33.5,28) {}; \node[tk] at (28.5,28) {}; \node[tk] at (23.5,28) {};
  \node[tk] at (33.5,14) {}; \node[tk] at (28.5,14) {};
  \node[tk] at (33.5,0) {};
  \node[anchor=west] at (44.5,28) {CPU};
  \node[anchor=west] at (44.5,14) {CPU};
  \node[anchor=west] at (44.5,0) {CPU};
  \node[desc] at (-3,28) {\iflangja{タイムスライス 8\,ms\\最も I/O 集約的\\最高優先度}{timeslice 8\,ms\\most I/O-intensive\\highest priority}};
  \node[desc] at (-3,14) {\iflangja{タイムスライス 16\,ms\\I/O と CPU の混合}{timeslice 16\,ms\\mix of I/O and CPU}};
  \node[desc] at (-3,0)  {\iflangja{タイムスライス $\infty$（FCFS）\\CPU 集約的\\最低優先度}{timeslice $\infty$ (FCFS)\\CPU-intensive\\lowest priority}};
  % new tasks enter the top level
  \draw[->] (2,40) node[above, font=\scriptsize\sffamily] {\iflangja{新しいタスク}{new task}} -- (2,31.5);
  % feedback
  \draw[->, thick] (56,27) to[bend left=40] (56,15);
  \draw[->, thick] (56,13) to[bend left=40] (56,1);
  \draw[->, thick, densely dashed] (66,1) to[bend right=40] (66,13);
  \draw[->, thick, densely dashed] (66,15) to[bend right=40] (66,27);
  % legend
  \begin{scope}[yshift=-12mm, font=\scriptsize\sffamily]
    \draw[->, thick] (0,0) -- (7,0);
    \node[anchor=west] at (8,0) {\iflangja{タイムスライスを使い切る：下の段へ}{uses its whole timeslice: moves down}};
    \draw[->, thick, densely dashed] (0,-5) -- (7,-5);
    \node[anchor=west] at (8,-5) {\iflangja{I/O のため早く CPU を手放す：上の段へ}{releases the CPU early for I/O: moves up}};
  \end{scope}
\end{tikzpicture}
   (width ~104 mm)
\begin{tikzpicture}[x=1mm, y=1mm, font=\footnotesize\sffamily,
  >={Stealth[length=1.6mm]},
  desc/.style={anchor=east, align=right, font=\scriptsize\sffamily},
  q/.style={draw, fill=ln-box},
  tk/.style={draw, circle, fill=white, minimum size=3.2mm, inner sep=0pt}]
  \foreach \y in {28,14,0} {
    \draw[q] (0,\y-3) rectangle (36,\y+3);
    \foreach \i in {1,...,4} \draw (36-5*\i,\y-3) -- ++(0,6);
    \draw[->] (36,\y) -- (44,\y);
  }
  \node[tk] at (33.5,28) {}; \node[tk] at (28.5,28) {}; \node[tk] at (23.5,28) {};
  \node[tk] at (33.5,14) {}; \node[tk] at (28.5,14) {};
  \node[tk] at (33.5,0) {};
  \node[anchor=west] at (44.5,28) {CPU};
  \node[anchor=west] at (44.5,14) {CPU};
  \node[anchor=west] at (44.5,0) {CPU};
  \node[desc] at (-3,28) {\iflangja{タイムスライス 8\,ms\\最も I/O 集約的\\最高優先度}{timeslice 8\,ms\\most I/O-intensive\\highest priority}};
  \node[desc] at (-3,14) {\iflangja{タイムスライス 16\,ms\\I/O と CPU の混合}{timeslice 16\,ms\\mix of I/O and CPU}};
  \node[desc] at (-3,0)  {\iflangja{タイムスライス $\infty$（FCFS）\\CPU 集約的\\最低優先度}{timeslice $\infty$ (FCFS)\\CPU-intensive\\lowest priority}};
  % new tasks enter the top level
  \draw[->] (2,40) node[above, font=\scriptsize\sffamily] {\iflangja{新しいタスク}{new task}} -- (2,31.5);
  % feedback
  \draw[->, thick] (56,27) to[bend left=40] (56,15);
  \draw[->, thick] (56,13) to[bend left=40] (56,1);
  \draw[->, thick, densely dashed] (66,1) to[bend right=40] (66,13);
  \draw[->, thick, densely dashed] (66,15) to[bend right=40] (66,27);
  % legend
  \begin{scope}[yshift=-12mm, font=\scriptsize\sffamily]
    \draw[->, thick] (0,0) -- (7,0);
    \node[anchor=west] at (8,0) {\iflangja{タイムスライスを使い切る：下の段へ}{uses its whole timeslice: moves down}};
    \draw[->, thick, densely dashed] (0,-5) -- (7,-5);
    \node[anchor=west] at (8,-5) {\iflangja{I/O のため早く CPU を手放す：上の段へ}{releases the CPU early for I/O: moves up}};
  \end{scope}
\end{tikzpicture}
   (width ~104 mm)
\begin{tikzpicture}[x=1mm, y=1mm, font=\footnotesize\sffamily,
  >={Stealth[length=1.6mm]},
  desc/.style={anchor=east, align=right, font=\scriptsize\sffamily},
  q/.style={draw, fill=ln-box},
  tk/.style={draw, circle, fill=white, minimum size=3.2mm, inner sep=0pt}]
  \foreach \y in {28,14,0} {
    \draw[q] (0,\y-3) rectangle (36,\y+3);
    \foreach \i in {1,...,4} \draw (36-5*\i,\y-3) -- ++(0,6);
    \draw[->] (36,\y) -- (44,\y);
  }
  \node[tk] at (33.5,28) {}; \node[tk] at (28.5,28) {}; \node[tk] at (23.5,28) {};
  \node[tk] at (33.5,14) {}; \node[tk] at (28.5,14) {};
  \node[tk] at (33.5,0) {};
  \node[anchor=west] at (44.5,28) {CPU};
  \node[anchor=west] at (44.5,14) {CPU};
  \node[anchor=west] at (44.5,0) {CPU};
  \node[desc] at (-3,28) {\iflangja{タイムスライス 8\,ms\\最も I/O 集約的\\最高優先度}{timeslice 8\,ms\\most I/O-intensive\\highest priority}};
  \node[desc] at (-3,14) {\iflangja{タイムスライス 16\,ms\\I/O と CPU の混合}{timeslice 16\,ms\\mix of I/O and CPU}};
  \node[desc] at (-3,0)  {\iflangja{タイムスライス $\infty$（FCFS）\\CPU 集約的\\最低優先度}{timeslice $\infty$ (FCFS)\\CPU-intensive\\lowest priority}};
  % new tasks enter the top level
  \draw[->] (2,40) node[above, font=\scriptsize\sffamily] {\iflangja{新しいタスク}{new task}} -- (2,31.5);
  % feedback
  \draw[->, thick] (56,27) to[bend left=40] (56,15);
  \draw[->, thick] (56,13) to[bend left=40] (56,1);
  \draw[->, thick, densely dashed] (66,1) to[bend right=40] (66,13);
  \draw[->, thick, densely dashed] (66,15) to[bend right=40] (66,27);
  % legend
  \begin{scope}[yshift=-12mm, font=\scriptsize\sffamily]
    \draw[->, thick] (0,0) -- (7,0);
    \node[anchor=west] at (8,0) {\iflangja{タイムスライスを使い切る：下の段へ}{uses its whole timeslice: moves down}};
    \draw[->, thick, densely dashed] (0,-5) -- (7,-5);
    \node[anchor=west] at (8,-5) {\iflangja{I/O のため早く CPU を手放す：上の段へ}{releases the CPU early for I/O: moves up}};
  \end{scope}
\end{tikzpicture}

  \caption{A multi-level feedback queue with three levels.  Each level
    serves its tasks with its own timeslice; tasks move between levels
    according to how they use their timeslices.}
  \label{fig:l07-mlfq}
\end{figure}

This design raises the question of how the scheduler knows whether a task
is I/O-intensive or CPU-intensive.  Heuristics based on the history of a
task do not help for a new task, nor for a task that changes its behaviour.
The solution is to let the scheduler learn the type of a task from the way
the task uses its timeslices:\caution{OSTEP's MLFQ differs in two details:
every level runs its tasks round robin, the lowest one included, and after a
period $S$ all tasks are moved back to the top level, so that none of them
starves (OSTEP ch.~8).}
\begin{enumerate}
  \item a new task enters the topmost queue;
  \item if the task yields the CPU before its timeslice expires, the choice
    of level was good, and the task stays at this level;
  \item if the task uses its entire timeslice, it is more CPU-intensive than
    the level assumes, and it is moved down one level;
  \item a task in a lower queue that releases the CPU early because it
    waits for I/O receives a priority boost and is moved up.
\end{enumerate}

\begin{definition}[Multi-level feedback queue]
  A \term{multi-level feedback queue}[多段フィードバックキュー]%
  \index{MLFQ|see{multi-level feedback queue}} (MLFQ) is a runqueue of
  several levels, each with its own scheduling policy, in which tasks move
  between the levels according to their observed use of the CPU.
\end{definition}

An MLFQ is not merely a set of priority queues.  Its levels treat tasks
differently, with timeslices of different lengths and FCFS at the lowest
level, and the feedback mechanism adjusts the level of each task over time.
The design goes back to the CTSS timesharing system of Fernando
Corbató \cite{corbato1962}, who received the Turing Award in 1990.

\begin{example}
  A new task enters the \SI{8}{\milli\second} level and computes for
  \SI{30}{\milli\second} without blocking.  It uses its whole first
  timeslice and moves to the \SI{16}{\milli\second} level, uses that
  timeslice as well and moves to the lowest level, where it completes its
  remaining \SI{6}{\milli\second}.  A task that computes for
  \SI{3}{\milli\second} between I/O operations never uses up an
  \SI{8}{\milli\second} timeslice and stays at the top level.
\end{example}

\Needspace{7\baselineskip}
\section{The Linux O(1) scheduler}
\label{sec:l07-o1}

The \term[O(1) scheduler]{$O(1)$ scheduler}[O(1)スケジューラ] of Linux
selects a task to run and adds a task to its runqueue in constant time,
regardless of the number of tasks in the system.\source{P3L1, Linux O(1)
scheduler; Bovet and Cesati \cite{bovet2005} ch.~7.}  It is a preemptive,
priority-based scheduler with 140 priority levels, numbered so that a
smaller number means a higher priority.  Levels 0 to 99 are for real-time
tasks, and levels 100 to 139 for timesharing tasks.

A user process has the default priority 120.  It can be adjusted through
the process's \term{nice value}[nice値], which ranges from $-20$ to 19 and
is added to the default, so that the priorities of timesharing tasks range
from 100 to 139.  A process can lower its own priority, for example with
the \texttt{nice} system call; raising it requires privileges.

The timeslice of a task depends on its priority: tasks of low priority
receive the shortest timeslices, and tasks of high priority the longest.
The scheduler also uses feedback, based on the \emph{sleep time} of a task,
the time it spends waiting or idling:
\begin{itemize}
  \item a task that sleeps for long periods is interactive, and its priority
    is boosted by up to 5 levels (its number decreases by up to 5);
  \item a task that sleeps little is compute-intensive, and its priority is
    lowered by up to 5 levels (its number increases by up to 5).
\end{itemize}

\begin{example}
  In the Linux~2.6 kernels described by Bovet and Cesati, a timesharing
  task whose static priority, the priority given by its nice value, is $p$
  receives a timeslice of $(140 - p) \times
  \SI{20}{\milli\second}$ if $p < 120$ and $(140 - p) \times
  \SI{5}{\milli\second}$ otherwise.  A process with nice value 5 has
  priority 125 and a timeslice of $15 \times 5 = \SI{75}{\milli\second}$; a
  process with nice value 0, \SI{100}{\milli\second}; a process with nice
  value $-10$, $30 \times 20 = \SI{600}{\milli\second}$.  If the process with
  nice value 5 is interactive, the feedback lets it run at a priority as high
  as 120; if it computes without sleeping, at a priority as low as
  130.\caution{The two numberings run in opposite directions: a larger
  priority number, and hence a larger nice value, means a lower priority,
  while a larger \texttt{sched\_priority} (1--99), the value used for
  \texttt{SCHED\_FIFO} and \texttt{SCHED\_RR}, means a higher one.}
\end{example}

\begin{figure}[htbp]
  \centering
  % Runqueue of the Linux O(1) scheduler: an active and an expired array of
% 140 priority lists, each with a bitmap of the non-empty lists.
% Usage: % Runqueue of the Linux O(1) scheduler: an active and an expired array of
% 140 priority lists, each with a bitmap of the non-empty lists.
% Usage: % Runqueue of the Linux O(1) scheduler: an active and an expired array of
% 140 priority lists, each with a bitmap of the non-empty lists.
% Usage: \input{figures/l07-o1}   (width ~112 mm)
\begin{tikzpicture}[x=1mm, y=1mm, font=\scriptsize\sffamily,
  >={Stealth[length=1.5mm]},
  bit/.style={draw, minimum width=4mm, minimum height=4mm, inner sep=0pt},
  slot/.style={draw, fill=ln-box, minimum width=5mm, minimum height=4mm, inner sep=0pt},
  tk/.style={draw, circle, fill=ln-accent!22, minimum size=3.6mm, inner sep=0pt},
  ttl/.style={font=\footnotesize\sffamily\bfseries}]
  % rows: label / y
  \def\rows{0/0, {\vdots}/-4.5, 99/-9, 100/-13.5, {\vdots}/-18, 120/-22.5, {\vdots}/-27, 139/-31.5}
  % ---------------- active array at x0 = 18
  \node[ttl, anchor=south] at (32,5) {\iflangja{アクティブ配列}{active array}};
  \node[anchor=south, text=ln-muted] at (18,2) {\iflangja{ビットマップ}{bitmap}};
  \foreach \l/\y in \rows { \node[anchor=east] at (29,\y) {\l}; }
  \foreach \y/\b in {0/0, -9/0, -13.5/1, -22.5/1, -31.5/1} {
    \node[bit] at (18,\y) {\b};
    \node[slot] at (33,\y) {};
  }
  \node[tk] (a1) at (41,-13.5) {};
  \draw[->] (33,-13.5) -- (a1);
  \node[tk] (a2) at (41,-22.5) {};
  \node[tk] (a3) at (48,-22.5) {};
  \draw[->] (33,-22.5) -- (a2); \draw[->] (a2) -- (a3);
  \node[tk] (a4) at (41,-31.5) {};
  \draw[->] (33,-31.5) -- (a4);
  % priority classes
  \draw (13,2) -- (11.5,2) -- (11.5,-11) -- (13,-11);
  \node[anchor=east, align=right] at (9.5,-4.5) {\iflangja{リアル\\タイム}{real-\\time}};
  \draw (13,-11.8) -- (11.5,-11.8) -- (11.5,-33.5) -- (13,-33.5);
  \node[anchor=east, align=right] at (9.5,-22.5) {\iflangja{タイム\\シェアリング}{time-\\sharing}};
  % ---------------- expired array at x0 = 76
  \node[ttl, anchor=south] at (90,5) {\iflangja{期限切れ配列}{expired array}};
  \node[anchor=south, text=ln-muted] at (76,2) {\iflangja{ビットマップ}{bitmap}};
  \foreach \l/\y in \rows { \node[anchor=east] at (87,\y) {\l}; }
  \foreach \y/\b in {0/0, -9/0, -13.5/0, -22.5/1, -31.5/0} {
    \node[bit] at (76,\y) {\b};
    \node[slot] at (91,\y) {};
  }
  \node[tk] (e1) at (99,-22.5) {};
  \draw[->] (91,-22.5) -- (e1);
  % a task whose timeslice expires moves to the expired array
  \draw[->, thick, ln-accent, rounded corners=2mm] (a3.south) -- (48,-37.5) -- (99,-37.5) -- (e1.south);
  \node[above, text=ln-accent, inner sep=1pt] at (62,-37.5) {\iflangja{タイムスライス満了}{timeslice expires}};
  % swap
  \draw[<->, thick] (44,-44) -- node[below, align=center] {\iflangja{アクティブ配列が空になったら交換}{swapped when the active array is empty}} (88,-44);
\end{tikzpicture}
   (width ~112 mm)
\begin{tikzpicture}[x=1mm, y=1mm, font=\scriptsize\sffamily,
  >={Stealth[length=1.5mm]},
  bit/.style={draw, minimum width=4mm, minimum height=4mm, inner sep=0pt},
  slot/.style={draw, fill=ln-box, minimum width=5mm, minimum height=4mm, inner sep=0pt},
  tk/.style={draw, circle, fill=ln-accent!22, minimum size=3.6mm, inner sep=0pt},
  ttl/.style={font=\footnotesize\sffamily\bfseries}]
  % rows: label / y
  \def\rows{0/0, {\vdots}/-4.5, 99/-9, 100/-13.5, {\vdots}/-18, 120/-22.5, {\vdots}/-27, 139/-31.5}
  % ---------------- active array at x0 = 18
  \node[ttl, anchor=south] at (32,5) {\iflangja{アクティブ配列}{active array}};
  \node[anchor=south, text=ln-muted] at (18,2) {\iflangja{ビットマップ}{bitmap}};
  \foreach \l/\y in \rows { \node[anchor=east] at (29,\y) {\l}; }
  \foreach \y/\b in {0/0, -9/0, -13.5/1, -22.5/1, -31.5/1} {
    \node[bit] at (18,\y) {\b};
    \node[slot] at (33,\y) {};
  }
  \node[tk] (a1) at (41,-13.5) {};
  \draw[->] (33,-13.5) -- (a1);
  \node[tk] (a2) at (41,-22.5) {};
  \node[tk] (a3) at (48,-22.5) {};
  \draw[->] (33,-22.5) -- (a2); \draw[->] (a2) -- (a3);
  \node[tk] (a4) at (41,-31.5) {};
  \draw[->] (33,-31.5) -- (a4);
  % priority classes
  \draw (13,2) -- (11.5,2) -- (11.5,-11) -- (13,-11);
  \node[anchor=east, align=right] at (9.5,-4.5) {\iflangja{リアル\\タイム}{real-\\time}};
  \draw (13,-11.8) -- (11.5,-11.8) -- (11.5,-33.5) -- (13,-33.5);
  \node[anchor=east, align=right] at (9.5,-22.5) {\iflangja{タイム\\シェアリング}{time-\\sharing}};
  % ---------------- expired array at x0 = 76
  \node[ttl, anchor=south] at (90,5) {\iflangja{期限切れ配列}{expired array}};
  \node[anchor=south, text=ln-muted] at (76,2) {\iflangja{ビットマップ}{bitmap}};
  \foreach \l/\y in \rows { \node[anchor=east] at (87,\y) {\l}; }
  \foreach \y/\b in {0/0, -9/0, -13.5/0, -22.5/1, -31.5/0} {
    \node[bit] at (76,\y) {\b};
    \node[slot] at (91,\y) {};
  }
  \node[tk] (e1) at (99,-22.5) {};
  \draw[->] (91,-22.5) -- (e1);
  % a task whose timeslice expires moves to the expired array
  \draw[->, thick, ln-accent, rounded corners=2mm] (a3.south) -- (48,-37.5) -- (99,-37.5) -- (e1.south);
  \node[above, text=ln-accent, inner sep=1pt] at (62,-37.5) {\iflangja{タイムスライス満了}{timeslice expires}};
  % swap
  \draw[<->, thick] (44,-44) -- node[below, align=center] {\iflangja{アクティブ配列が空になったら交換}{swapped when the active array is empty}} (88,-44);
\end{tikzpicture}
   (width ~112 mm)
\begin{tikzpicture}[x=1mm, y=1mm, font=\scriptsize\sffamily,
  >={Stealth[length=1.5mm]},
  bit/.style={draw, minimum width=4mm, minimum height=4mm, inner sep=0pt},
  slot/.style={draw, fill=ln-box, minimum width=5mm, minimum height=4mm, inner sep=0pt},
  tk/.style={draw, circle, fill=ln-accent!22, minimum size=3.6mm, inner sep=0pt},
  ttl/.style={font=\footnotesize\sffamily\bfseries}]
  % rows: label / y
  \def\rows{0/0, {\vdots}/-4.5, 99/-9, 100/-13.5, {\vdots}/-18, 120/-22.5, {\vdots}/-27, 139/-31.5}
  % ---------------- active array at x0 = 18
  \node[ttl, anchor=south] at (32,5) {\iflangja{アクティブ配列}{active array}};
  \node[anchor=south, text=ln-muted] at (18,2) {\iflangja{ビットマップ}{bitmap}};
  \foreach \l/\y in \rows { \node[anchor=east] at (29,\y) {\l}; }
  \foreach \y/\b in {0/0, -9/0, -13.5/1, -22.5/1, -31.5/1} {
    \node[bit] at (18,\y) {\b};
    \node[slot] at (33,\y) {};
  }
  \node[tk] (a1) at (41,-13.5) {};
  \draw[->] (33,-13.5) -- (a1);
  \node[tk] (a2) at (41,-22.5) {};
  \node[tk] (a3) at (48,-22.5) {};
  \draw[->] (33,-22.5) -- (a2); \draw[->] (a2) -- (a3);
  \node[tk] (a4) at (41,-31.5) {};
  \draw[->] (33,-31.5) -- (a4);
  % priority classes
  \draw (13,2) -- (11.5,2) -- (11.5,-11) -- (13,-11);
  \node[anchor=east, align=right] at (9.5,-4.5) {\iflangja{リアル\\タイム}{real-\\time}};
  \draw (13,-11.8) -- (11.5,-11.8) -- (11.5,-33.5) -- (13,-33.5);
  \node[anchor=east, align=right] at (9.5,-22.5) {\iflangja{タイム\\シェアリング}{time-\\sharing}};
  % ---------------- expired array at x0 = 76
  \node[ttl, anchor=south] at (90,5) {\iflangja{期限切れ配列}{expired array}};
  \node[anchor=south, text=ln-muted] at (76,2) {\iflangja{ビットマップ}{bitmap}};
  \foreach \l/\y in \rows { \node[anchor=east] at (87,\y) {\l}; }
  \foreach \y/\b in {0/0, -9/0, -13.5/0, -22.5/1, -31.5/0} {
    \node[bit] at (76,\y) {\b};
    \node[slot] at (91,\y) {};
  }
  \node[tk] (e1) at (99,-22.5) {};
  \draw[->] (91,-22.5) -- (e1);
  % a task whose timeslice expires moves to the expired array
  \draw[->, thick, ln-accent, rounded corners=2mm] (a3.south) -- (48,-37.5) -- (99,-37.5) -- (e1.south);
  \node[above, text=ln-accent, inner sep=1pt] at (62,-37.5) {\iflangja{タイムスライス満了}{timeslice expires}};
  % swap
  \draw[<->, thick] (44,-44) -- node[below, align=center] {\iflangja{アクティブ配列が空になったら交換}{swapped when the active array is empty}} (88,-44);
\end{tikzpicture}

  \caption{The runqueue of the $O(1)$ scheduler.  Each array has one list
    per priority level and a bitmap of the non-empty lists.  A task whose
    timeslice expires moves to the expired array; when the active array is
    empty, the two arrays are swapped.}
  \label{fig:l07-o1}
\end{figure}

The runqueue consists of two arrays of task lists, with one list for each
priority level (\cref{fig:l07-o1}).  The \term{active array}[アクティブ配列]
is the one from which the scheduler picks the next task.  A bitmap records
which of its 140 lists are non-empty; the scheduler finds the first set bit
and takes the first task of that list, and adding a task appends it to a
list and sets a bit.  Both operations take constant time, because the
number of levels is fixed.  A task remains in the active array until its
timeslice expires: a task that blocks before then returns to the active
array, with the rest of its timeslice, when it becomes ready again.

The \term{expired array}[期限切れ配列] holds the tasks whose timeslices
have expired and that are therefore inactive.  When the timeslice of a task
expires, the scheduler recomputes its priority and timeslice and moves it to
the expired array.  When no more tasks remain in the active array, the
scheduler swaps the two arrays, which requires only exchanging two
pointers.  This arrangement also gives tasks of low priority a chance to
run: a task of high priority, although it receives a long timeslice, moves
to the expired array once it has used that timeslice, and waits there
until every task in the active array, including those of low priority, has
used its own.\margin{In the Linux~2.6 kernels a task judged interactive
was normally returned to the active array instead, unless the tasks in the
expired array had already waited too long.}

The $O(1)$ scheduler was introduced by Ingo Molnar in the Linux~2.5
development series and became the scheduler of Linux~2.6.  As workloads
became more interactive, with applications such as streaming media, its
jitter became a problem, and in Linux~2.6.23 it was replaced by the
Completely Fair Scheduler, also written by Molnar.

\Needspace{7\baselineskip}
\section{The Linux Completely Fair Scheduler}
\label{sec:l07-cfs}

The $O(1)$ scheduler has two problems.\source{P3L1, Linux CFS scheduler;
OSTEP ch.~9.}  First, the performance of interactive tasks is poor: once a
task is in the expired array, it is not scheduled until all tasks in the
active array have used their timeslices, so an interactive task may wait an
unpredictable time before it runs.  Second, it makes no guarantee of
fairness, in the sense that within a given time interval every task should
run for an amount of time proportional to its priority.

The \term{Completely Fair Scheduler}[CFS]%
\index{CFS|see{Completely Fair Scheduler}} (CFS) addresses both problems.
It was the default Linux scheduler for non-real-time tasks from version
2.6.23 to 6.5.  Its runqueue is a \term{red-black tree}[赤黒木], a
self-balancing binary search tree in which no path from the root to a leaf is more than
twice as long as any other, so that all paths have approximately the same
length.  The tree is ordered by the virtual runtime of the tasks
(\cref{fig:l07-cfs}).

\begin{definition}[Virtual runtime]
  The \term{virtual runtime}[仮想実行時間] (\texttt{vruntime}) of a task is
  the time the task has spent running on the CPU, measured with nanosecond
  granularity and scaled according to the priority of the task.
\end{definition}

\begin{figure}[htbp]
  \centering
  % CFS runqueue: a red-black tree of the ready tasks ordered by virtual
% runtime (ms).
% Usage: % CFS runqueue: a red-black tree of the ready tasks ordered by virtual
% runtime (ms).
% Usage: % CFS runqueue: a red-black tree of the ready tasks ordered by virtual
% runtime (ms).
% Usage: \input{figures/l07-cfs}   (width ~108 mm)
\begin{tikzpicture}[x=1mm, y=1mm, font=\footnotesize\sffamily,
  >={Stealth[length=1.6mm]},
  bn/.style={draw, circle, fill=ln-muted, text=white, minimum size=7mm, inner sep=0pt},
  rn/.style={draw, circle, fill=ln-caution!80, text=white, minimum size=7mm, inner sep=0pt},
  lab/.style={font=\scriptsize\sffamily, align=center}]
  \node[bn] (n24) at (48,0)   {24};
  \node[bn] (n15) at (28,-13) {15};
  \node[bn] (n38) at (68,-13) {38};
  \node[rn] (n9)  at (18,-26) {9};
  \node[rn] (n19) at (38,-26) {19};
  \node[rn] (n30) at (58,-26) {30};
  \node[rn] (n45) at (78,-26) {45};
  \draw (n24) -- (n15); \draw (n24) -- (n38);
  \draw (n15) -- (n9);  \draw (n15) -- (n19);
  \draw (n38) -- (n30); \draw (n38) -- (n45);
  % leftmost node
  \draw[->, thick, ln-accent] (n9.west) -- ++(-9,0) node[left, lab, text=ln-accent]
    {\iflangja{最左の節点：\\vruntime 最小\\→ 次に実行}{leftmost node:\\smallest vruntime\\$\to$ runs next}};
  % a preempted task, whose vruntime has just passed that of the leftmost
  % node, is reinserted
  \node[bn, dashed, fill=white, text=black, draw=black] (p) at (98,-8) {12};
  \node[lab, anchor=south] at (98,-3.5) {\iflangja{プリエンプト\\されたタスク}{preempted\\task}};
  \node[draw, dashed, circle, minimum size=7mm, inner sep=0pt] (ins) at (25,-38) {};
  \draw[dashed] (n9) -- (ins);
  \draw[->, thick, densely dashed] (p.south) .. controls (98,-48) and (58,-48) .. (ins.east);
  \node[lab] at (72,-38) {\iflangja{挿入}{insert}};
  \node[lab, text=ln-muted, anchor=west] at (0,-48) {\iflangja{節点の値：仮想実行時間（ms）}{node values: virtual runtime (ms)}};
\end{tikzpicture}
   (width ~108 mm)
\begin{tikzpicture}[x=1mm, y=1mm, font=\footnotesize\sffamily,
  >={Stealth[length=1.6mm]},
  bn/.style={draw, circle, fill=ln-muted, text=white, minimum size=7mm, inner sep=0pt},
  rn/.style={draw, circle, fill=ln-caution!80, text=white, minimum size=7mm, inner sep=0pt},
  lab/.style={font=\scriptsize\sffamily, align=center}]
  \node[bn] (n24) at (48,0)   {24};
  \node[bn] (n15) at (28,-13) {15};
  \node[bn] (n38) at (68,-13) {38};
  \node[rn] (n9)  at (18,-26) {9};
  \node[rn] (n19) at (38,-26) {19};
  \node[rn] (n30) at (58,-26) {30};
  \node[rn] (n45) at (78,-26) {45};
  \draw (n24) -- (n15); \draw (n24) -- (n38);
  \draw (n15) -- (n9);  \draw (n15) -- (n19);
  \draw (n38) -- (n30); \draw (n38) -- (n45);
  % leftmost node
  \draw[->, thick, ln-accent] (n9.west) -- ++(-9,0) node[left, lab, text=ln-accent]
    {\iflangja{最左の節点：\\vruntime 最小\\→ 次に実行}{leftmost node:\\smallest vruntime\\$\to$ runs next}};
  % a preempted task, whose vruntime has just passed that of the leftmost
  % node, is reinserted
  \node[bn, dashed, fill=white, text=black, draw=black] (p) at (98,-8) {12};
  \node[lab, anchor=south] at (98,-3.5) {\iflangja{プリエンプト\\されたタスク}{preempted\\task}};
  \node[draw, dashed, circle, minimum size=7mm, inner sep=0pt] (ins) at (25,-38) {};
  \draw[dashed] (n9) -- (ins);
  \draw[->, thick, densely dashed] (p.south) .. controls (98,-48) and (58,-48) .. (ins.east);
  \node[lab] at (72,-38) {\iflangja{挿入}{insert}};
  \node[lab, text=ln-muted, anchor=west] at (0,-48) {\iflangja{節点の値：仮想実行時間（ms）}{node values: virtual runtime (ms)}};
\end{tikzpicture}
   (width ~108 mm)
\begin{tikzpicture}[x=1mm, y=1mm, font=\footnotesize\sffamily,
  >={Stealth[length=1.6mm]},
  bn/.style={draw, circle, fill=ln-muted, text=white, minimum size=7mm, inner sep=0pt},
  rn/.style={draw, circle, fill=ln-caution!80, text=white, minimum size=7mm, inner sep=0pt},
  lab/.style={font=\scriptsize\sffamily, align=center}]
  \node[bn] (n24) at (48,0)   {24};
  \node[bn] (n15) at (28,-13) {15};
  \node[bn] (n38) at (68,-13) {38};
  \node[rn] (n9)  at (18,-26) {9};
  \node[rn] (n19) at (38,-26) {19};
  \node[rn] (n30) at (58,-26) {30};
  \node[rn] (n45) at (78,-26) {45};
  \draw (n24) -- (n15); \draw (n24) -- (n38);
  \draw (n15) -- (n9);  \draw (n15) -- (n19);
  \draw (n38) -- (n30); \draw (n38) -- (n45);
  % leftmost node
  \draw[->, thick, ln-accent] (n9.west) -- ++(-9,0) node[left, lab, text=ln-accent]
    {\iflangja{最左の節点：\\vruntime 最小\\→ 次に実行}{leftmost node:\\smallest vruntime\\$\to$ runs next}};
  % a preempted task, whose vruntime has just passed that of the leftmost
  % node, is reinserted
  \node[bn, dashed, fill=white, text=black, draw=black] (p) at (98,-8) {12};
  \node[lab, anchor=south] at (98,-3.5) {\iflangja{プリエンプト\\されたタスク}{preempted\\task}};
  \node[draw, dashed, circle, minimum size=7mm, inner sep=0pt] (ins) at (25,-38) {};
  \draw[dashed] (n9) -- (ins);
  \draw[->, thick, densely dashed] (p.south) .. controls (98,-48) and (58,-48) .. (ins.east);
  \node[lab] at (72,-38) {\iflangja{挿入}{insert}};
  \node[lab, text=ln-muted, anchor=west] at (0,-48) {\iflangja{節点の値：仮想実行時間（ms）}{node values: virtual runtime (ms)}};
\end{tikzpicture}

  \caption{The runqueue of CFS, a red-black tree of the ready tasks ordered
    by virtual runtime (grey nodes are black, red nodes red).  The leftmost
    task runs next; a preempted task is inserted at the position given by
    its virtual runtime.}
  \label{fig:l07-cfs}
\end{figure}

The algorithm of CFS is simple:\margin{Since Linux~6.6 the fair
scheduling class selects tasks by the EEVDF algorithm, which retains the
virtual runtime accounting of CFS.}
\begin{itemize}
  \item it always schedules the task of the leftmost node, which has the
    smallest virtual runtime, that is, the task that has so far received
    the least CPU time;
  \item while the task runs, CFS periodically increases its virtual runtime
    and compares it with the virtual runtime of the leftmost task in the
    tree;
  \item if the running task's virtual runtime is still smaller, the task
    continues to run; if it is larger, the task is preempted and inserted
    into the tree at the position its virtual runtime dictates, and the
    leftmost task runs.
\end{itemize}

Priorities are expressed by the rate at which virtual runtime progresses.
For tasks of low priority, virtual runtime advances faster than real time,
so they soon lose the CPU; for tasks of high priority it advances more
slowly, so they run longer.  A single tree thus holds the tasks of all
priorities.  Selecting a task takes constant time, since the scheduler keeps
a pointer to the leftmost node, and adding a task takes $O(\log N)$ time for
$N$ tasks.

\begin{example}
  CFS derives the rate from a weight for each nice value: nice value 0 has
  weight 1024, and each step of the nice value changes the weight by a
  factor of about 1.25; nice value 5 has weight 335.\tip{The factor is
  chosen so that one nice level is worth about \SI{10}{\percent} of the CPU:
  one level apart gives $1.25/2.25 \approx 55:45$, five levels apart the
  75:25 here.}  A task of weight $w$
  that runs for a time $\Delta t$ advances its virtual runtime by
  $\Delta t \cdot 1024/w$.  If task A (nice 0) and task B (nice 5) each run
  for \SI{6}{\milli\second}, the virtual runtime of A grows by
  \SI{6}{\milli\second} but that of B by $6 \cdot 1024/335 \approx
  \SI{18.3}{\milli\second}$.  Because CFS keeps the virtual runtimes of the
  tasks close, B runs about a third as long as A, and the two receive
  $1024/1359 \approx \SI{75}{\percent}$ and \SI{25}{\percent} of the CPU.
\end{example}

\begin{table}[htbp]
  \centering
  \caption{The two Linux schedulers.}
  \label{tab:l07-linux}
  \small
  \begin{tabular}{@{}>{\raggedright\arraybackslash}p{0.22\linewidth}>{\raggedright\arraybackslash}p{0.34\linewidth}>{\raggedright\arraybackslash}p{0.34\linewidth}@{}}
    \toprule
    & $O(1)$ scheduler & CFS \\
    \midrule
    Linux versions & 2.6.0--2.6.22 & 2.6.23--6.5 \\
    Runqueue & active and expired arrays of 140 lists & red-black tree ordered
      by virtual runtime \\
    Select a task & $O(1)$ & $O(1)$ \\
    Add a task & $O(1)$ & $O(\log N)$ \\
    Priorities & selection order, that is, the highest non-empty list, and
      timeslice length; sleep-time feedback adjusts the priority & rate of
      virtual runtime \\
    Interactive tasks & may wait for the whole active array & run as soon
      as they have the least virtual runtime \\
    \bottomrule
  \end{tabular}
\end{table}

\Needspace{7\baselineskip}
\section{Scheduling on multiprocessors}
\label{sec:l07-multi}

Scheduling becomes more complex when there are several
CPUs.\source{P3L1, scheduling on multiprocessors; OSTEP ch.~10;
Silberschatz et al.\ ch.~5.}  Two architectures must be distinguished
(\cref{fig:l07-multicpu}).

\begin{definition}[Shared memory multiprocessor]
  A \term{shared memory multiprocessor}[共有メモリ型マルチプロセッサ] has
  several CPUs, each with its own private caches, all of which share main
  memory.
\end{definition}

\begin{definition}[Multicore processor]
  A \term{multicore processor}[マルチコアプロセッサ] contains several
  cores, each of which is a CPU with its own private caches; the cores
  share the last-level cache (LLC) of the processor and main memory.
\end{definition}

\begin{figure}[htbp]
  \centering
  % (a) Shared memory multiprocessor: CPUs with private caches share memory.
% (b) Multicore processor: cores with private caches share the LLC.
% Usage: % (a) Shared memory multiprocessor: CPUs with private caches share memory.
% (b) Multicore processor: cores with private caches share the LLC.
% Usage: % (a) Shared memory multiprocessor: CPUs with private caches share memory.
% (b) Multicore processor: cores with private caches share the LLC.
% Usage: \input{figures/l07-multicpu}   (width ~112 mm)
\begin{tikzpicture}[x=1mm, y=1mm, font=\footnotesize\sffamily,
  cpu/.style={draw, fill=ln-accent!22, minimum width=17mm, minimum height=6mm, inner sep=1pt},
  cache/.style={draw, fill=ln-box, minimum width=17mm, minimum height=4.5mm,
                inner sep=1pt, font=\scriptsize\sffamily},
  mem/.style={draw, fill=ln-box, minimum height=7mm, inner sep=1pt},
  lab/.style={font=\scriptsize\sffamily, align=center}]
  % ---------------- (a) SMP
  \foreach \x/\n in {12/0, 40/1} {
    \node[cpu] (a\n) at (\x,0) {CPU\,\n};
    \node[cache, anchor=north] (al\n) at (a\n.south) {L1, L2};
    \node[cache, anchor=north] (ac\n) at (al\n.south) {LLC};
    \draw (ac\n.south) -- (\x,-20);
  }
  \draw[line width=1.2pt] (4,-20) -- (48,-20);
  \node[mem, minimum width=36mm] (am) at (26,-28) {\iflangja{メモリ（DRAM）}{memory (DRAM)}};
  \draw (am.north) -- (26,-20);
  \node[lab] at (26,-38) {\iflangja{(a) 共有メモリ型マルチプロセッサ}{(a) shared memory multiprocessor}};
  % ---------------- (b) multicore
  \draw[rounded corners=2pt, dashed, ln-muted] (64,4.5) rectangle (112,-17);
  \node[lab, anchor=north east, text=ln-muted] at (111.5,-17.5) {\iflangja{プロセッサ}{processor}};
  \foreach \x/\n in {76/0, 100/1} {
    \node[cpu] (b\n) at (\x,0) {\iflangja{コア}{core}\,\n};
    \node[cache, anchor=north] (bl\n) at (b\n.south) {L1, L2};
  }
  \node[cache, minimum width=41mm] (llc) at (88,-13) {\iflangja{LLC（共有）}{LLC (shared)}};
  \draw (bl0.south) -- (bl0.south |- llc.north);
  \draw (bl1.south) -- (bl1.south |- llc.north);
  \node[mem, minimum width=36mm] (bm) at (88,-28) {\iflangja{メモリ（DRAM）}{memory (DRAM)}};
  \draw (llc.south) -- (bm.north);
  \node[lab] at (88,-38) {\iflangja{(b) マルチコアプロセッサ}{(b) multicore processor}};
\end{tikzpicture}
   (width ~112 mm)
\begin{tikzpicture}[x=1mm, y=1mm, font=\footnotesize\sffamily,
  cpu/.style={draw, fill=ln-accent!22, minimum width=17mm, minimum height=6mm, inner sep=1pt},
  cache/.style={draw, fill=ln-box, minimum width=17mm, minimum height=4.5mm,
                inner sep=1pt, font=\scriptsize\sffamily},
  mem/.style={draw, fill=ln-box, minimum height=7mm, inner sep=1pt},
  lab/.style={font=\scriptsize\sffamily, align=center}]
  % ---------------- (a) SMP
  \foreach \x/\n in {12/0, 40/1} {
    \node[cpu] (a\n) at (\x,0) {CPU\,\n};
    \node[cache, anchor=north] (al\n) at (a\n.south) {L1, L2};
    \node[cache, anchor=north] (ac\n) at (al\n.south) {LLC};
    \draw (ac\n.south) -- (\x,-20);
  }
  \draw[line width=1.2pt] (4,-20) -- (48,-20);
  \node[mem, minimum width=36mm] (am) at (26,-28) {\iflangja{メモリ（DRAM）}{memory (DRAM)}};
  \draw (am.north) -- (26,-20);
  \node[lab] at (26,-38) {\iflangja{(a) 共有メモリ型マルチプロセッサ}{(a) shared memory multiprocessor}};
  % ---------------- (b) multicore
  \draw[rounded corners=2pt, dashed, ln-muted] (64,4.5) rectangle (112,-17);
  \node[lab, anchor=north east, text=ln-muted] at (111.5,-17.5) {\iflangja{プロセッサ}{processor}};
  \foreach \x/\n in {76/0, 100/1} {
    \node[cpu] (b\n) at (\x,0) {\iflangja{コア}{core}\,\n};
    \node[cache, anchor=north] (bl\n) at (b\n.south) {L1, L2};
  }
  \node[cache, minimum width=41mm] (llc) at (88,-13) {\iflangja{LLC（共有）}{LLC (shared)}};
  \draw (bl0.south) -- (bl0.south |- llc.north);
  \draw (bl1.south) -- (bl1.south |- llc.north);
  \node[mem, minimum width=36mm] (bm) at (88,-28) {\iflangja{メモリ（DRAM）}{memory (DRAM)}};
  \draw (llc.south) -- (bm.north);
  \node[lab] at (88,-38) {\iflangja{(b) マルチコアプロセッサ}{(b) multicore processor}};
\end{tikzpicture}
   (width ~112 mm)
\begin{tikzpicture}[x=1mm, y=1mm, font=\footnotesize\sffamily,
  cpu/.style={draw, fill=ln-accent!22, minimum width=17mm, minimum height=6mm, inner sep=1pt},
  cache/.style={draw, fill=ln-box, minimum width=17mm, minimum height=4.5mm,
                inner sep=1pt, font=\scriptsize\sffamily},
  mem/.style={draw, fill=ln-box, minimum height=7mm, inner sep=1pt},
  lab/.style={font=\scriptsize\sffamily, align=center}]
  % ---------------- (a) SMP
  \foreach \x/\n in {12/0, 40/1} {
    \node[cpu] (a\n) at (\x,0) {CPU\,\n};
    \node[cache, anchor=north] (al\n) at (a\n.south) {L1, L2};
    \node[cache, anchor=north] (ac\n) at (al\n.south) {LLC};
    \draw (ac\n.south) -- (\x,-20);
  }
  \draw[line width=1.2pt] (4,-20) -- (48,-20);
  \node[mem, minimum width=36mm] (am) at (26,-28) {\iflangja{メモリ（DRAM）}{memory (DRAM)}};
  \draw (am.north) -- (26,-20);
  \node[lab] at (26,-38) {\iflangja{(a) 共有メモリ型マルチプロセッサ}{(a) shared memory multiprocessor}};
  % ---------------- (b) multicore
  \draw[rounded corners=2pt, dashed, ln-muted] (64,4.5) rectangle (112,-17);
  \node[lab, anchor=north east, text=ln-muted] at (111.5,-17.5) {\iflangja{プロセッサ}{processor}};
  \foreach \x/\n in {76/0, 100/1} {
    \node[cpu] (b\n) at (\x,0) {\iflangja{コア}{core}\,\n};
    \node[cache, anchor=north] (bl\n) at (b\n.south) {L1, L2};
  }
  \node[cache, minimum width=41mm] (llc) at (88,-13) {\iflangja{LLC（共有）}{LLC (shared)}};
  \draw (bl0.south) -- (bl0.south |- llc.north);
  \draw (bl1.south) -- (bl1.south |- llc.north);
  \node[mem, minimum width=36mm] (bm) at (88,-28) {\iflangja{メモリ（DRAM）}{memory (DRAM)}};
  \draw (llc.south) -- (bm.north);
  \node[lab] at (88,-38) {\iflangja{(b) マルチコアプロセッサ}{(b) multicore processor}};
\end{tikzpicture}

  \caption{(a) A shared memory multiprocessor: every CPU has its own
    caches.  (b) A multicore processor: the cores have private caches and
    share the last-level cache.}
  \label{fig:l07-multicpu}
\end{figure}

To the operating system, each core of a multicore processor is a CPU, and a
system may combine both architectures, with several multicore processors.

\subsection{Cache affinity and load balancing}

The performance of a task depends on the state of the caches.  While a task
runs on a CPU, the CPU's caches fill with the task's data.  If the task
runs on the same CPU the next time, it finds a hot cache (Lesson~2); if it
runs on a different CPU, it starts with a cold cache and runs more slowly
until that cache has been filled.\margin{A task that ran within the last
\texttt{sched\_migration\_cost\_ns}, \SI{500}{\micro\second} by default,
counts as cache-hot, and the load balancer leaves it where it is.}

\begin{definition}[Cache affinity]
  The benefit that a task gains from running on the CPU whose caches hold
  its data is its \term{cache affinity}[キャッシュアフィニティ].  A scheduler
  preserves it by keeping each task on the same CPU as much as possible.
\end{definition}

Schedulers for multiprocessors therefore use a hierarchical architecture.
Each CPU has its own runqueue, and a scheduler for that CPU selects tasks
only from it, so that a task keeps running on the same CPU.  Above the
per-CPU schedulers, a \term{load balancing}[負荷分散] component moves tasks
between the runqueues so that no CPU is overloaded while another has too
little to do.  It acts on the lengths of the runqueues, or when a CPU becomes
idle and takes tasks from the runqueues of other CPUs.\margin{On Linux an
application can restrict a thread to a set of CPUs with
\texttt{sched\_setaffinity}.}

\subsection{NUMA-aware scheduling}

A multiprocessor may also have several memory nodes (\cref{fig:l07-numa}).

\begin{definition}[NUMA]
  In a system with \term{non-uniform memory access}[NUMA]%
  \index{NUMA|see{non-uniform memory access}} (NUMA) main memory consists
  of several memory nodes, each of which is close to a socket of one or more
  processors.  Access from a CPU to its local memory node is faster than
  access to a remote node, which passes through an interconnect between the
  sockets, such as Intel's QuickPath Interconnect.
\end{definition}

\begin{figure}[htbp]
  \centering
  % Hierarchical scheduling on a NUMA system: per-CPU runqueues, a load
% balancer, and two memory nodes connected by an interconnect.
% Usage: % Hierarchical scheduling on a NUMA system: per-CPU runqueues, a load
% balancer, and two memory nodes connected by an interconnect.
% Usage: % Hierarchical scheduling on a NUMA system: per-CPU runqueues, a load
% balancer, and two memory nodes connected by an interconnect.
% Usage: \input{figures/l07-numa}   (width ~112 mm)
\begin{tikzpicture}[x=1mm, y=1mm, font=\footnotesize\sffamily,
  >={Stealth[length=1.6mm]},
  cpu/.style={draw, fill=ln-accent!22, minimum width=16mm, minimum height=6mm, inner sep=1pt},
  mem/.style={draw, fill=ln-box, minimum width=38mm, minimum height=7mm, inner sep=1pt},
  nd/.style={draw, dashed, rounded corners=2pt, ln-muted},
  lab/.style={font=\scriptsize\sffamily, align=center},
  tk/.style={draw, circle, fill=white, minimum size=2.6mm, inner sep=0pt}]
  % load balancer
  \draw[fill=ln-box] (0,36) rectangle (112,42);
  \node at (56,39) {\iflangja{負荷分散}{load balancing}};
  % per-CPU runqueues and CPUs
  \foreach \x/\n/\k in {13/0/3, 36/1/1, 76/2/2, 99/3/0} {
    \draw[<->] (\x,35.5) -- (\x,29.5);
    \draw[fill=white] (\x-8,24) rectangle (\x+8,29);
    \foreach \i in {1,...,3} \draw (\x+8-4*\i,24) -- ++(0,5);
    \ifnum\k>0 \foreach \i in {1,...,\k} \node[tk] at (\x+10-4*\i,26.5) {}; \fi
    \draw[->] (\x,24) -- (\x,17);
    \node[cpu] (c\n) at (\x,14) {CPU\,\n};
  }
  \node[lab] at (56,26.5) {\iflangja{CPU ごとの\\ランキュー}{per-CPU\\runqueues}};
  % memory nodes
  \node[mem] (m0) at (24.5,-6) {\iflangja{メモリノード 0}{memory node 0}};
  \node[mem] (m1) at (87.5,-6) {\iflangja{メモリノード 1}{memory node 1}};
  \draw[nd] (2,20) rectangle (47,-12);
  \draw[nd] (65,20) rectangle (110,-12);
  \node[lab, anchor=north, text=ln-muted] at (24.5,-12.5) {\iflangja{ノード 0（ソケット 0）}{node 0 (socket 0)}};
  \node[lab, anchor=north, text=ln-muted] at (87.5,-12.5) {\iflangja{ノード 1（ソケット 1）}{node 1 (socket 1)}};
  % local access
  \draw[<->, thick] (c0.south) -- node[lab, right] {\iflangja{ローカル：\\速い}{local:\\fast}} (c0.south |- m0.north);
  \draw (c1.south) -- (c1.south |- m0.north);
  \draw (c2.south) -- (c2.south |- m1.north);
  \draw (c3.south) -- (c3.south |- m1.north);
  % interconnect
  \draw[line width=2pt, ln-muted] (m0.east) -- (m1.west);
  \node[lab, anchor=north, text=ln-muted] at (56,-12.5) {\iflangja{インターコネクト}{interconnect}};
  \draw[ln-muted, densely dotted] (56,-6) -- (56,-12.3);
  % remote access
  \draw[<->, thick, densely dashed, ln-caution] (c2.south west) -- (m0.north east);
  \node[lab, text=ln-caution] at (53.5,10.5) {\iflangja{リモート：\\遅い}{remote:\\slow}};
\end{tikzpicture}
   (width ~112 mm)
\begin{tikzpicture}[x=1mm, y=1mm, font=\footnotesize\sffamily,
  >={Stealth[length=1.6mm]},
  cpu/.style={draw, fill=ln-accent!22, minimum width=16mm, minimum height=6mm, inner sep=1pt},
  mem/.style={draw, fill=ln-box, minimum width=38mm, minimum height=7mm, inner sep=1pt},
  nd/.style={draw, dashed, rounded corners=2pt, ln-muted},
  lab/.style={font=\scriptsize\sffamily, align=center},
  tk/.style={draw, circle, fill=white, minimum size=2.6mm, inner sep=0pt}]
  % load balancer
  \draw[fill=ln-box] (0,36) rectangle (112,42);
  \node at (56,39) {\iflangja{負荷分散}{load balancing}};
  % per-CPU runqueues and CPUs
  \foreach \x/\n/\k in {13/0/3, 36/1/1, 76/2/2, 99/3/0} {
    \draw[<->] (\x,35.5) -- (\x,29.5);
    \draw[fill=white] (\x-8,24) rectangle (\x+8,29);
    \foreach \i in {1,...,3} \draw (\x+8-4*\i,24) -- ++(0,5);
    \ifnum\k>0 \foreach \i in {1,...,\k} \node[tk] at (\x+10-4*\i,26.5) {}; \fi
    \draw[->] (\x,24) -- (\x,17);
    \node[cpu] (c\n) at (\x,14) {CPU\,\n};
  }
  \node[lab] at (56,26.5) {\iflangja{CPU ごとの\\ランキュー}{per-CPU\\runqueues}};
  % memory nodes
  \node[mem] (m0) at (24.5,-6) {\iflangja{メモリノード 0}{memory node 0}};
  \node[mem] (m1) at (87.5,-6) {\iflangja{メモリノード 1}{memory node 1}};
  \draw[nd] (2,20) rectangle (47,-12);
  \draw[nd] (65,20) rectangle (110,-12);
  \node[lab, anchor=north, text=ln-muted] at (24.5,-12.5) {\iflangja{ノード 0（ソケット 0）}{node 0 (socket 0)}};
  \node[lab, anchor=north, text=ln-muted] at (87.5,-12.5) {\iflangja{ノード 1（ソケット 1）}{node 1 (socket 1)}};
  % local access
  \draw[<->, thick] (c0.south) -- node[lab, right] {\iflangja{ローカル：\\速い}{local:\\fast}} (c0.south |- m0.north);
  \draw (c1.south) -- (c1.south |- m0.north);
  \draw (c2.south) -- (c2.south |- m1.north);
  \draw (c3.south) -- (c3.south |- m1.north);
  % interconnect
  \draw[line width=2pt, ln-muted] (m0.east) -- (m1.west);
  \node[lab, anchor=north, text=ln-muted] at (56,-12.5) {\iflangja{インターコネクト}{interconnect}};
  \draw[ln-muted, densely dotted] (56,-6) -- (56,-12.3);
  % remote access
  \draw[<->, thick, densely dashed, ln-caution] (c2.south west) -- (m0.north east);
  \node[lab, text=ln-caution] at (53.5,10.5) {\iflangja{リモート：\\遅い}{remote:\\slow}};
\end{tikzpicture}
   (width ~112 mm)
\begin{tikzpicture}[x=1mm, y=1mm, font=\footnotesize\sffamily,
  >={Stealth[length=1.6mm]},
  cpu/.style={draw, fill=ln-accent!22, minimum width=16mm, minimum height=6mm, inner sep=1pt},
  mem/.style={draw, fill=ln-box, minimum width=38mm, minimum height=7mm, inner sep=1pt},
  nd/.style={draw, dashed, rounded corners=2pt, ln-muted},
  lab/.style={font=\scriptsize\sffamily, align=center},
  tk/.style={draw, circle, fill=white, minimum size=2.6mm, inner sep=0pt}]
  % load balancer
  \draw[fill=ln-box] (0,36) rectangle (112,42);
  \node at (56,39) {\iflangja{負荷分散}{load balancing}};
  % per-CPU runqueues and CPUs
  \foreach \x/\n/\k in {13/0/3, 36/1/1, 76/2/2, 99/3/0} {
    \draw[<->] (\x,35.5) -- (\x,29.5);
    \draw[fill=white] (\x-8,24) rectangle (\x+8,29);
    \foreach \i in {1,...,3} \draw (\x+8-4*\i,24) -- ++(0,5);
    \ifnum\k>0 \foreach \i in {1,...,\k} \node[tk] at (\x+10-4*\i,26.5) {}; \fi
    \draw[->] (\x,24) -- (\x,17);
    \node[cpu] (c\n) at (\x,14) {CPU\,\n};
  }
  \node[lab] at (56,26.5) {\iflangja{CPU ごとの\\ランキュー}{per-CPU\\runqueues}};
  % memory nodes
  \node[mem] (m0) at (24.5,-6) {\iflangja{メモリノード 0}{memory node 0}};
  \node[mem] (m1) at (87.5,-6) {\iflangja{メモリノード 1}{memory node 1}};
  \draw[nd] (2,20) rectangle (47,-12);
  \draw[nd] (65,20) rectangle (110,-12);
  \node[lab, anchor=north, text=ln-muted] at (24.5,-12.5) {\iflangja{ノード 0（ソケット 0）}{node 0 (socket 0)}};
  \node[lab, anchor=north, text=ln-muted] at (87.5,-12.5) {\iflangja{ノード 1（ソケット 1）}{node 1 (socket 1)}};
  % local access
  \draw[<->, thick] (c0.south) -- node[lab, right] {\iflangja{ローカル：\\速い}{local:\\fast}} (c0.south |- m0.north);
  \draw (c1.south) -- (c1.south |- m0.north);
  \draw (c2.south) -- (c2.south |- m1.north);
  \draw (c3.south) -- (c3.south |- m1.north);
  % interconnect
  \draw[line width=2pt, ln-muted] (m0.east) -- (m1.west);
  \node[lab, anchor=north, text=ln-muted] at (56,-12.5) {\iflangja{インターコネクト}{interconnect}};
  \draw[ln-muted, densely dotted] (56,-6) -- (56,-12.3);
  % remote access
  \draw[<->, thick, densely dashed, ln-caution] (c2.south west) -- (m0.north east);
  \node[lab, text=ln-caution] at (53.5,10.5) {\iflangja{リモート：\\遅い}{remote:\\slow}};
\end{tikzpicture}

  \caption{Hierarchical scheduling on a NUMA system.  Every CPU schedules
    tasks from its own runqueue, and load balancing moves tasks between the
    runqueues.  A CPU reaches the memory of its own node directly and the
    memory of the other node through the interconnect.}
  \label{fig:l07-numa}
\end{figure}

With \term{NUMA-aware scheduling}[NUMA対応スケジューリング] the scheduler
keeps a task on a CPU close to the memory node that holds the task's state,
so that its memory accesses remain local.  Load balancing must take this into account: moving a
task to a CPU on another node may balance the runqueues but slows down
every access of the task to main memory, which then crosses the
interconnect.

\begin{example}
  A task spends \SI{40}{\percent} of its execution time accessing main
  memory, and
  a remote access takes 1.5 times as long as a local one.  If the load
  balancer moves the task to a CPU of the other node while its memory stays
  where it is, the task's execution time grows by a factor of
  $0.6 + 0.4 \times 1.5 = 1.2$, that is, by \SI{20}{\percent}.
\end{example}

\Needspace{5\baselineskip}
\section{Hyperthreading}
\label{sec:l07-smt}

The CPUs must be multiplexed among tasks because the execution context of a
task, its register values such as the program counter and the stack
pointer, is held in the registers of a CPU while the task
runs.\source{P3L1, hyperthreading; Silberschatz et al.\ ch.~5; Fedorova et
al.\ \cite{fedorova2004}.}  To switch to another task, the context must be
saved and the other context loaded.  Hardware can instead provide several
sets of registers, and thus several hardware-supported execution contexts,
on one CPU.  The CPU remains a single CPU, with one set of execution
resources shared by the contexts, but passing from one context to another
needs no saving and restoring of state.

\begin{definition}[Simultaneous multithreading]
  In \term{simultaneous multithreading}[同時マルチスレッディング]%
  \index{SMT|see{simultaneous multithreading}}%
  \index{hyperthreading|see{simultaneous multithreading}} (SMT) a CPU
  provides several hardware execution contexts, or hardware threads, each
  with its own register state, which share the CPU's pipeline and caches.
\end{definition}

The technique is also called hardware multithreading, hyperthreading, the
name used by Intel, or chip multithreading.  More precisely,
\term{chip multithreading}[チップマルチスレッディング]%
\index{CMT|see{chip multithreading}} (CMT) refers to processors that combine
several cores with several hardware threads per core, the platforms studied
by Fedorova et al.; in practice the names are often used interchangeably.  Processors
with up to eight hardware threads per core have been built, but two is the
most common configuration, and SMT can be enabled or disabled at boot
time.\margin{Intel's first implementation cost less than \SI{5}{\percent} of
the die area for about \SI{25}{\percent} more throughput (Marr et al.,
\emph{Intel Technology Journal}, 2002).}

Designs differ in how closely the contexts share the pipeline.  The simplest
interleave them, giving the pipeline to another context as often as every
cycle; in simultaneous multithreading proper, such as Intel's
Hyper-Threading, instructions of several hardware threads occupy the stages
of the pipeline in the same cycle.\margin{The cores simulated by Fedorova et
al.\ interleave four hardware threads cycle by cycle.}

Each hardware thread appears to the operating system as a separate virtual
CPU, on which the scheduler can place a task.  The scheduler must therefore
decide which tasks to co-schedule on the hardware threads of the same CPU.
A context switch between hardware threads takes on the order of a few
cycles, whereas loading data from memory after a cache miss takes on the
order of 100 cycles.  By the argument of Lesson~3, switching pays off when
the idle time exceeds twice the cost of a switch (\cref{eq:l03-idle}); it
pays off by a wide margin here.  With SMT, a CPU can therefore hide the
latency of memory accesses, just as multithreading on a single CPU hides the
latency of I/O.

\Needspace{7\baselineskip}
\section{Scheduling for hyperthreaded platforms}
\label{sec:l07-smt-sched}

Which tasks should share a CPU?  A simple model of a multithreaded CPU, in
the spirit of the platforms studied by Fedorova et al., answers the
question.\source{P3L1, scheduling for hyperthreading platforms; Fedorova et
al.\ \cite{fedorova2004}.}  It rests on four assumptions.
\begin{enumerate}
  \item A thread issues an instruction on every CPU cycle, and the pipeline
    executes at most one instruction per cycle.
  \item A memory access of a memory-bound thread takes 4 cycles, during which
    the thread issues no instructions.  The small value keeps the timing
    diagram short; it is the delay of a load served by the L1 cache rather
    than the hundred or so cycles of an access that reaches memory.
  \item Switching between hardware threads is instantaneous.
  \item The CPU has two hardware threads.
\end{enumerate}
In this model, a CPU-bound task needs the pipeline on every cycle.  A
\term{memory-bound task}[メモリバウンドタスク] is one whose progress is
limited by memory accesses: it spends most of its cycles waiting for data to
arrive from memory.  \Cref{fig:l07-smt} shows the three possible pairings.

\begin{figure}[htbp]
  \centering
  % Co-scheduling two threads on the two hardware threads of an SMT CPU.
% One instruction per cycle on the shared pipeline; a memory access takes
% 4 cycles; switching between hardware threads is instantaneous.
% Usage: % Co-scheduling two threads on the two hardware threads of an SMT CPU.
% One instruction per cycle on the shared pipeline; a memory access takes
% 4 cycles; switching between hardware threads is instantaneous.
% Usage: % Co-scheduling two threads on the two hardware threads of an SMT CPU.
% One instruction per cycle on the shared pipeline; a memory access takes
% 4 cycles; switching between hardware threads is instantaneous.
% Usage: \input{figures/l07-smt}   (width ~104 mm)
\begin{tikzpicture}[x=8mm, y=1mm, font=\scriptsize\sffamily,
  ex/.style={draw, fill=ln-accent!30},
  mw/.style={draw, fill=ln-box},
  rd/.style={draw, fill=white},
  lab/.style={anchor=east, font=\scriptsize\sffamily},
  ttl/.style={anchor=south west, font=\footnotesize\sffamily\bfseries, inner sep=1pt}]
  % cycle numbers
  \node[lab] at (-0.1,60) {\iflangja{サイクル}{cycle}};
  \foreach \c in {1,...,10} \node at (\c-0.5,60) {\c};
  % cells: \cell{row y}{cycle}{style}{text}
  \def\cell#1#2#3#4{\draw[#3] (#2-1,#1) rectangle (#2,#1+5); \node at (#2-0.5,#1+2.5) {#4};}
  % ---------- (a) two CPU-bound threads
  \node[ttl] at (0,51.5) {\iflangja{(a) CPU バウンド 2 つ}{(a) two CPU-bound threads}};
  \node[lab] at (-0.1,48.5) {HT\,0: C};
  \node[lab] at (-0.1,42.5) {HT\,1: C};
  \foreach \c in {1,3,5,7,9} { \cell{46}{\c}{ex}{C} \cell{40}{\c}{rd}{} }
  \foreach \c in {2,4,6,8,10} { \cell{46}{\c}{rd}{} \cell{40}{\c}{ex}{C} }
  % ---------- (b) two memory-bound threads
  \node[ttl] at (0,32.5) {\iflangja{(b) メモリバウンド 2 つ}{(b) two memory-bound threads}};
  \node[lab] at (-0.1,29.5) {HT\,0: M};
  \node[lab] at (-0.1,23.5) {HT\,1: M};
  \foreach \c in {1,6} { \cell{27}{\c}{ex}{M} }
  \foreach \c in {2,3,4,5,7,8,9,10} { \cell{27}{\c}{mw}{} }
  \foreach \c in {2,7} { \cell{21}{\c}{ex}{M} }
  \foreach \c in {1,3,4,5,6,8,9,10} { \cell{21}{\c}{mw}{} }
  \foreach \c in {3,4,5,8,9,10} { \node[text=ln-caution] at (\c-0.5,18.6) {\iflangja{空き}{idle}}; }
  % ---------- (c) one CPU-bound and one memory-bound thread
  \node[ttl] at (0,9.5) {\iflangja{(c) CPU バウンドとメモリバウンド}{(c) one CPU-bound, one memory-bound thread}};
  \node[lab] at (-0.1,6.5) {HT\,0: M};
  \node[lab] at (-0.1,0.5) {HT\,1: C};
  \foreach \c in {1,6} { \cell{4}{\c}{ex}{M} }
  \foreach \c in {2,3,4,5,7,8,9,10} { \cell{4}{\c}{mw}{} }
  \foreach \c in {1,6} { \cell{-2}{\c}{rd}{} }
  \foreach \c in {2,3,4,5,7,8,9,10} { \cell{-2}{\c}{ex}{C} }
  % ---------- legend
  \begin{scope}[yshift=-9mm]
    \draw[ex] (0,-1.5) rectangle (0.5,1.5);
    \node[anchor=west] at (0.55,0) {\iflangja{命令を実行}{executes an instruction}};
    \draw[mw] (3.9,-1.5) rectangle (4.4,1.5);
    \node[anchor=west] at (4.45,0) {\iflangja{メモリを待つ}{waits for memory}};
    \draw[rd] (7.2,-1.5) rectangle (7.7,1.5);
    \node[anchor=west] at (7.75,0) {\iflangja{実行可能だが待つ}{ready, not executing}};
  \end{scope}
\end{tikzpicture}
   (width ~104 mm)
\begin{tikzpicture}[x=8mm, y=1mm, font=\scriptsize\sffamily,
  ex/.style={draw, fill=ln-accent!30},
  mw/.style={draw, fill=ln-box},
  rd/.style={draw, fill=white},
  lab/.style={anchor=east, font=\scriptsize\sffamily},
  ttl/.style={anchor=south west, font=\footnotesize\sffamily\bfseries, inner sep=1pt}]
  % cycle numbers
  \node[lab] at (-0.1,60) {\iflangja{サイクル}{cycle}};
  \foreach \c in {1,...,10} \node at (\c-0.5,60) {\c};
  % cells: \cell{row y}{cycle}{style}{text}
  \def\cell#1#2#3#4{\draw[#3] (#2-1,#1) rectangle (#2,#1+5); \node at (#2-0.5,#1+2.5) {#4};}
  % ---------- (a) two CPU-bound threads
  \node[ttl] at (0,51.5) {\iflangja{(a) CPU バウンド 2 つ}{(a) two CPU-bound threads}};
  \node[lab] at (-0.1,48.5) {HT\,0: C};
  \node[lab] at (-0.1,42.5) {HT\,1: C};
  \foreach \c in {1,3,5,7,9} { \cell{46}{\c}{ex}{C} \cell{40}{\c}{rd}{} }
  \foreach \c in {2,4,6,8,10} { \cell{46}{\c}{rd}{} \cell{40}{\c}{ex}{C} }
  % ---------- (b) two memory-bound threads
  \node[ttl] at (0,32.5) {\iflangja{(b) メモリバウンド 2 つ}{(b) two memory-bound threads}};
  \node[lab] at (-0.1,29.5) {HT\,0: M};
  \node[lab] at (-0.1,23.5) {HT\,1: M};
  \foreach \c in {1,6} { \cell{27}{\c}{ex}{M} }
  \foreach \c in {2,3,4,5,7,8,9,10} { \cell{27}{\c}{mw}{} }
  \foreach \c in {2,7} { \cell{21}{\c}{ex}{M} }
  \foreach \c in {1,3,4,5,6,8,9,10} { \cell{21}{\c}{mw}{} }
  \foreach \c in {3,4,5,8,9,10} { \node[text=ln-caution] at (\c-0.5,18.6) {\iflangja{空き}{idle}}; }
  % ---------- (c) one CPU-bound and one memory-bound thread
  \node[ttl] at (0,9.5) {\iflangja{(c) CPU バウンドとメモリバウンド}{(c) one CPU-bound, one memory-bound thread}};
  \node[lab] at (-0.1,6.5) {HT\,0: M};
  \node[lab] at (-0.1,0.5) {HT\,1: C};
  \foreach \c in {1,6} { \cell{4}{\c}{ex}{M} }
  \foreach \c in {2,3,4,5,7,8,9,10} { \cell{4}{\c}{mw}{} }
  \foreach \c in {1,6} { \cell{-2}{\c}{rd}{} }
  \foreach \c in {2,3,4,5,7,8,9,10} { \cell{-2}{\c}{ex}{C} }
  % ---------- legend
  \begin{scope}[yshift=-9mm]
    \draw[ex] (0,-1.5) rectangle (0.5,1.5);
    \node[anchor=west] at (0.55,0) {\iflangja{命令を実行}{executes an instruction}};
    \draw[mw] (3.9,-1.5) rectangle (4.4,1.5);
    \node[anchor=west] at (4.45,0) {\iflangja{メモリを待つ}{waits for memory}};
    \draw[rd] (7.2,-1.5) rectangle (7.7,1.5);
    \node[anchor=west] at (7.75,0) {\iflangja{実行可能だが待つ}{ready, not executing}};
  \end{scope}
\end{tikzpicture}
   (width ~104 mm)
\begin{tikzpicture}[x=8mm, y=1mm, font=\scriptsize\sffamily,
  ex/.style={draw, fill=ln-accent!30},
  mw/.style={draw, fill=ln-box},
  rd/.style={draw, fill=white},
  lab/.style={anchor=east, font=\scriptsize\sffamily},
  ttl/.style={anchor=south west, font=\footnotesize\sffamily\bfseries, inner sep=1pt}]
  % cycle numbers
  \node[lab] at (-0.1,60) {\iflangja{サイクル}{cycle}};
  \foreach \c in {1,...,10} \node at (\c-0.5,60) {\c};
  % cells: \cell{row y}{cycle}{style}{text}
  \def\cell#1#2#3#4{\draw[#3] (#2-1,#1) rectangle (#2,#1+5); \node at (#2-0.5,#1+2.5) {#4};}
  % ---------- (a) two CPU-bound threads
  \node[ttl] at (0,51.5) {\iflangja{(a) CPU バウンド 2 つ}{(a) two CPU-bound threads}};
  \node[lab] at (-0.1,48.5) {HT\,0: C};
  \node[lab] at (-0.1,42.5) {HT\,1: C};
  \foreach \c in {1,3,5,7,9} { \cell{46}{\c}{ex}{C} \cell{40}{\c}{rd}{} }
  \foreach \c in {2,4,6,8,10} { \cell{46}{\c}{rd}{} \cell{40}{\c}{ex}{C} }
  % ---------- (b) two memory-bound threads
  \node[ttl] at (0,32.5) {\iflangja{(b) メモリバウンド 2 つ}{(b) two memory-bound threads}};
  \node[lab] at (-0.1,29.5) {HT\,0: M};
  \node[lab] at (-0.1,23.5) {HT\,1: M};
  \foreach \c in {1,6} { \cell{27}{\c}{ex}{M} }
  \foreach \c in {2,3,4,5,7,8,9,10} { \cell{27}{\c}{mw}{} }
  \foreach \c in {2,7} { \cell{21}{\c}{ex}{M} }
  \foreach \c in {1,3,4,5,6,8,9,10} { \cell{21}{\c}{mw}{} }
  \foreach \c in {3,4,5,8,9,10} { \node[text=ln-caution] at (\c-0.5,18.6) {\iflangja{空き}{idle}}; }
  % ---------- (c) one CPU-bound and one memory-bound thread
  \node[ttl] at (0,9.5) {\iflangja{(c) CPU バウンドとメモリバウンド}{(c) one CPU-bound, one memory-bound thread}};
  \node[lab] at (-0.1,6.5) {HT\,0: M};
  \node[lab] at (-0.1,0.5) {HT\,1: C};
  \foreach \c in {1,6} { \cell{4}{\c}{ex}{M} }
  \foreach \c in {2,3,4,5,7,8,9,10} { \cell{4}{\c}{mw}{} }
  \foreach \c in {1,6} { \cell{-2}{\c}{rd}{} }
  \foreach \c in {2,3,4,5,7,8,9,10} { \cell{-2}{\c}{ex}{C} }
  % ---------- legend
  \begin{scope}[yshift=-9mm]
    \draw[ex] (0,-1.5) rectangle (0.5,1.5);
    \node[anchor=west] at (0.55,0) {\iflangja{命令を実行}{executes an instruction}};
    \draw[mw] (3.9,-1.5) rectangle (4.4,1.5);
    \node[anchor=west] at (4.45,0) {\iflangja{メモリを待つ}{waits for memory}};
    \draw[rd] (7.2,-1.5) rectangle (7.7,1.5);
    \node[anchor=west] at (7.75,0) {\iflangja{実行可能だが待つ}{ready, not executing}};
  \end{scope}
\end{tikzpicture}

  \caption{Co-scheduling two threads on the hardware threads (HT) of one
    SMT CPU, with C for a CPU-bound and M for a memory-bound thread.  Each
    column is a cycle of the shared pipeline.}
  \label{fig:l07-smt}
\end{figure}

\begin{description}
  \item[Two CPU-bound threads.]  Both threads are ready on every cycle, but
    the pipeline executes only one instruction per cycle.  The threads
    interfere, contending for the pipeline: each progresses at half the
    speed it would have alone, a performance degradation by a factor of 2.
    The memory component meanwhile is idle.
  \item[Two memory-bound threads.]  Both threads spend most of their cycles
    waiting for memory, and during 6 of the 10 cycles in the figure neither
    has an instruction to execute.  These CPU cycles are wasted.
  \item[A CPU-bound and a memory-bound thread.]  The CPU-bound thread uses
    the cycles in which the memory-bound thread waits: it executes on 8 of
    10 cycles, while the memory-bound thread progresses as fast as it would
    alone.  No cycle is idle, and both the CPU and the memory are well
    utilized.  Some interference remains, since the CPU-bound thread loses
    one cycle in five, but it is small.
\end{description}

\Needspace{6\baselineskip}
\begin{keypoint}
  On an SMT CPU, co-scheduling a CPU-bound with a memory-bound task avoids
  both contention for the pipeline and idle cycles.  The scheduler should
  therefore mix the two kinds of tasks on the hardware threads of each CPU.
\end{keypoint}

\Needspace{7\baselineskip}
\section{Hardware performance counters}
\label{sec:l07-counters}

Mixing CPU-bound and memory-bound tasks requires the scheduler to know
which is which.\source{P3L1, CPU bound or memory bound, scheduling with
hardware counters; Fedorova et al.}  The history-based heuristics used for
I/O-bound tasks do not work.  Sleep time, as used by the $O(1)$ scheduler,
does not reveal a memory-bound task, because a thread that waits for memory
is not sleeping: it waits inside the processor, invisible to the operating
system.  Nor can software keep track of individual memory accesses, which
would take far too much time.  The information must come from the hardware.

\begin{definition}[Hardware performance counter]
  A \term{hardware performance counter}[ハードウェア性能カウンタ] is a
  register of the processor that is updated as the processor executes, and
  counts events relevant to performance.
\end{definition}

Processors typically provide several counters, and the counters available
differ between architectures.  They count, for example:
\begin{itemize}
  \item cache misses at each level of the cache hierarchy, L1, L2 and so on
    up to the LLC;
  \item instructions and cycles, from which the instructions per cycle (IPC)
    or its inverse, the cycles per instruction, are derived;
  \item power and energy consumption, on some processors.
\end{itemize}
Software interfaces and tools make the counters accessible, for example
\texttt{oprofile}, whose documentation lists the counters available on
different architectures, and the Linux \texttt{perf} tool.\tip{The command
\texttt{perf stat -e cycles,instructions,LLC-load-misses} \emph{cmd} reports
these counts for a program \emph{cmd}.}

From counter values the scheduler can estimate what kind of resources a
task needs.  A task with few LLC misses finds its data in the cache and is
likely CPU-bound; a task with many LLC misses keeps waiting for memory and
is likely memory-bound.  A single counter is rarely conclusive.  Schedulers
therefore combine several counters in models with thresholds tuned for each
architecture, based on well-understood characteristics of workloads.

\begin{definition}[Cycles per instruction]
  The \term{cycles per instruction}[CPI]%
  \index{CPI|see{cycles per instruction}} (CPI) of a task is the number of CPU cycles it takes per
  instruction executed, the ratio of the cycle count to the instruction
  count.
\end{definition}

\begin{example}
  During an interval of \SI{100}{\milli\second}, the counters report that
  task A executed $4.2 \times 10^8$ instructions in $4.6 \times 10^8$
  cycles, with 0.2 LLC misses per 1000 instructions, and that task B
  executed $0.9 \times 10^8$ instructions in the same number of cycles, with
  25 LLC misses per 1000 instructions.  A has a CPI of about 1.1 and B of
  about 5.1.  Both indicators classify A as CPU-bound and B as memory-bound,
  and a scheduler following \cref{sec:l07-smt-sched} would place them on the
  two hardware threads of the same CPU.
\end{example}

\Needspace{7\baselineskip}
\section{Is cycles per instruction a useful metric?}
\label{sec:l07-cpi}

Fedorova et al.\ proposed to distinguish the tasks by a single
counter-derived metric, the CPI.\source{P3L1, CPI experiment, CPI
experiment results; Fedorova et al.}  The idea is attractive: a
memory-bound task spends many cycles waiting for memory and has a high CPI,
while a CPU-bound task has a low one: 1 in a model that executes one
instruction per cycle, and below 1 on a superscalar processor.  A scheduler could then mix
tasks with different CPIs on the hardware threads of each core.

To test the idea, they simulated a chip multithreading processor with four
cores of four hardware threads each, sixteen hardware threads in total.  The
workload was synthetic and consisted of threads with widely different CPIs.
The threads were assigned to the cores in several ways, ranging from
assignments in which every core ran a mix of CPIs to assignments in which
every core ran threads of the same CPI.  Performance was measured as the
aggregate instructions per cycle of the sixteen hardware threads, at most 4
on four cores.

The assignments with mixed CPIs achieved the highest IPC, and those in
which the threads on each core had similar CPIs the lowest.  With these
workloads, scheduling by CPI is clearly beneficial, and the experiment
confirms that contention for the processor pipeline matters.  Realistic
workloads, however, behave differently.  The CPIs measured for real
applications, such as those of standard benchmark suites, are not nearly as
distinct: they cluster in a narrow range.\sidenote{The workshop paper itself
was optimistic about CPI: it reports that measured CPI varies by an order of
magnitude between applications, and that the near-term CPI of a thread can
be predicted.  The limitations it names are different ones: CPI reveals
neither contention for the caches nor contention for the floating-point
unit, and it is distorted when measured on a busy system.}  CPI therefore
cannot tell the tasks of a real
workload apart well enough to guide the scheduler, and what works for the
synthetic workload does not work in practice.

\begin{note}
  The study nevertheless leaves three lessons.  Tasks on SMT processors
  contend for the processor pipeline; hardware counters can characterize a
  workload; and schedulers must be aware of resource contention, not merely
  balance the load.  Later work by the same group found the use of the
  last-level cache over time to be a better indicator of contention than
  CPI.
\end{note}

\begin{keypoint}[Lesson summary]
  A scheduler selects the next task from its runqueue, and the runqueue
  structure is designed together with the algorithm.  FCFS and SJF run
  tasks to completion; preemption and priorities serve urgent tasks but
  bring starvation, avoided by aging, and priority inversion, avoided by
  priority inheritance.  Round robin with timeslices makes a system
  responsive; CPU-bound tasks favour long timeslices and I/O-bound tasks
  short ones, which a multi-level feedback queue provides by learning the
  type of each task.  Linux replaced its $O(1)$ scheduler, with active and
  expired arrays, by CFS, which runs the task of least virtual runtime from a
  red-black tree.  On multiprocessors, per-CPU runqueues preserve cache
  affinity, load balancing spreads the tasks, and NUMA-aware scheduling keeps
  tasks near their memory.  On SMT processors, CPU-bound and memory-bound
  tasks should share a CPU; hardware performance counters help to tell them
  apart, although CPI alone is not a sufficient metric.
\end{keypoint}

\end{document}
